% !TeX program = lualatex
\documentclass[12pt, a4paper]{article}
\usepackage{fullpage}
\usepackage{subfiles}
\usepackage{fontspec}
\usepackage{libertine}
\usepackage{xcolor}
\usepackage{GotIn}
\usepackage{geometry}
\usepackage{multicol}
\usepackage{multicolrule}
\usepackage{graphicx}
\usepackage{enumitem}
\usepackage[autocompile]{gregoriotex}
\usepackage[latin,french]{babel}


\geometry{top=2cm, bottom=2cm}
% \pagestyle{empty}

\definecolor{red}{HTML}{C70039}
% \input GoudyIn.fd
% \newcommand*\initfamily{\usefont{U}{GoudyIn}{xl}{n}}

\input Acorn.fd
\newcommand*\initfamily{\usefont{U}{Acorn}{xl}{n}}
% cette ligne ajoute de l'espace entre les portées
% \grechangedim{baselineskip}{60pt}{scalable}

\begin{document}
  \gresetlinecolor{gregoriocolor}

  \begin{titlepage}\centering
    \vspace*{\fill}\
    \huge Secondes Vêpres\\
    \smallskip
    \begin{Large}
      \textit{
        des Dimanches de l'Avent\\
      }
    \end{Large}
    \medskip
    \large et\\
    \medskip
    \LARGE Salut du Saint-Sacrement\\
    \bigskip
    \vspace*{\fill}
    \begin{figure}[h!]
      \centering
      \includegraphics[width=7cm]{../epiphanie_septuagesime_careme/logo.png}
    \end{figure}
    \centering \normalsize Paroisse Saint Roch\\
    \bigskip
    \begin{Large}
      \centering Ne pas emporter
    \end{Large}
  \end{titlepage}

  \newpage
  \vspace*{\fill}
  \begin{center}
    \normalsize\textit{
      Livret latin-français
    }
  \end{center}
  \newpage

  \begin{center}
    \textcolor{red}{\large{Invitatoire.}}
  \end{center}

  % greillumination: remplace la première lettre, ici par une font ornementale
  \greillumination{\initfamily\fontsize{11mm}{11mm}\selectfont D}
  \gregorioscore{or--deus_in_adjutorium_(tonus_festivus)--solesmes_1961.1}

  \begin{center}
    \small{
    \emph{
      Dieu, venez à mon aide ; Seigneur, hatez-vous de me secourir.\\
      Gloire au Père, au Fils et au Saint Esprit, comme il était au commencement, maintenant et toujour et dans les siècles des siècles. Allelúia\\
    }
  }
  \end{center}

  \bigskip

  \begin{center}
    \textcolor{red}{\footnotesize\textit{
    Le célébrant entonne la première antienne. Les chantres entonnent les antiennes suivantes, ainsi que le premier verset de chaque psaumes.
  }}
  \end{center}

  \vspace*{\fill}\
  \begin{center}
    \greseparator{3}{30}
  \end{center}
  \vspace*{\fill}\

  \newpage
  \normalsize

  \begin{center}
    \begin{LARGE}
      I\textsuperscript{er} dimanche de l'Avent
    \end{LARGE}
  \end{center}
  \medskip

  % ===== DEBUT Antienne =========
  \greillumination{\initfamily\fontsize{11mm}{11mm}\selectfont I}
  \gregorioscore{antiennes/an--in_illa_die--solesmes_1961}
  \begin{center}
    \footnotesize{
      \textit{
        En ce jour-là, la douceur gouttera des montagnes, et le lait et le miel couleront des collines, alléluia. 
      }
    }
  \end{center}
  % \medskip
  % ===== FIN Antienne ===========

  % ===== DEBUT psaume ===========
  % gresetinitiallines : avec le parametre à 0, supprime l'ornement
  \gresetinitiallines{0}

  \begin{center}
    \normalsize{Psaume 109.}
  \end{center}

  \gregorioscore{psaumes/psaume109-VIIIG}
  \begin{enumerate}[label=\textcolor{red}{\emph{\arabic*}}]
    \setcounter{enumi}{2}
    \item Donec ponam inimícos \textbf{tu}os,~* scabéllum pe\textit{dum} \textit{tu}\textbf{ó}rum.

    \item Virgam virtútis tuæ emíttet Dóminus ex \textbf{Si}on:~* domináre in médio inimicó\textit{rum} \textit{tu}\textbf{ó}rum.
    
    \item Tecum princípium in die virtútis tuæ in splendóribus sanc\textbf{tó}rum:~* ex útero ante lucíferum \textit{gé}\textit{nu}\textbf{i} te.
    
    \item Jurávit Dóminus, et non pœnitébit \textbf{e}um:~* Tu es sacérdos in ætérnum secúndum órdi\textit{nem} \textit{Mel}\textbf{chí}sedech.
    
    \item Dóminus a dextris \textbf{tu}is,~* confrégit in die iræ \textit{su}\textit{æ} \textbf{re}ges.
    
    \item Judicábit in natiónibus, implébit ru\textbf{í}nas:~* conquassábit cápita in ter\textit{ra} \textit{mul}\textbf{tó}rum.
    
    \item De torrénte in via \textbf{bi}bet:~* proptérea exal\textit{tá}\textit{bit} \textbf{ca}put.
    
    \item Glória Patri, et \textbf{Fí}lio,~* et Spirí\textit{tu}\textit{i} \textbf{Sanc}to.
    
    \item Sicut erat in princípio, et nunc, et \textbf{sem}per,~* et in sǽcula sæcu\textit{ló}\textit{rum}. \textbf{A}men.
  \end{enumerate}
  \medskip
  \grecommentary{\textit{Reprise de l'Antienne.}}
  \gabcsnippet{(c4) In(g) il(ggf~)la(gh) di(h)e(g'_[oh:h]) (,) stil(g)lá(g)bunt(g') mon(g)tes(h_i) dul(g)cé(h)di(g)nem,(f.) (;) et(f') col(g)les(h') flu(j)ent(g') lac(h) et(g) mel,(gv_[oh:h]FD') al(f)le(fg)lú(g.){ia}.(g.) (::)}

  \newpage
  \vspace*{\fill}\
  \begin{normalsize}
    \begin{center}
      \par \textit{L'Éternel a dit à mon Seigneur: Assieds-toi à ma droite.}
      \medskip
      \begin{enumerate}[label=\textcolor{red}{\emph{\arabic*}}]
        \item \textit{Jusqu'à ce que je mette tes ennemis pour le marchepied de tes pieds.}
        \item \textit{L'Éternel enverra de Sion la verge de ta force: Domine au milieu de tes ennemis!}
        \item \textit{Ton peuple sera un peuple de franche volonté, au jour de ta puissance, en sainte magnificence. Du sein de l'aurore te viendra la rosée de ta jeunesse.}
        \item \textit{L'Éternel a juré, et il ne se repentira point: Tu es sacrificateur pour toujours, selon l'ordre de Melchisédec.}
        \item \textit{Le Seigneur, à ta droite, brisera les rois au jour de sa colère.}
        \item \textit{Il jugera parmi les nations, il remplira tout de corps morts, il brisera le chef d'un grand pays.}
        \item \textit{Il boira du torrent dans le chemin, c'est pourquoi il lèvera haut la tête.}
        \item \textit{Gloire au Père, au Fils, et au Saint Esprit, }
        \item \textit{Comme il était au commencement, maintenant et toujours, et dans les siècles des siècles. Amen. }
      \end{enumerate}
    \end{center}
  \end{normalsize}
  \vspace*{\fill}\
  \newpage

  % ===== DEBUT Antienne =========
  \gresetinitiallines{1}
  \greillumination{\initfamily\fontsize{11mm}{11mm}\selectfont J}
  \gregorioscore{antiennes/an--jucundare_filia--solesmes_1961}
  \begin{center}
    \footnotesize{
      \textit{
        Réjouissez-vous, fille de Sion ; et tressaillez de joie, fille de Jérusalem, alléluia. 
      }
    }
  \end{center}
  % \medskip
  % ===== FIN Antienne ===========

  % ===== DEBUT psaume ===========
  % gresetinitiallines : avec le parametre à 0, supprime l'ornement
  \gresetinitiallines{0}

  \begin{center}
    \normalsize{Psaume 110.}
  \end{center}

  \gregorioscore{psaumes/psaume110-VIIIG}

  \begin{enumerate}[label=\textcolor{red}{\emph{\arabic*}}]
    \setcounter{enumi}{1}
    \item Magna ópera \textbf{Dó}mini:~* exquisíta in omnes volun\textit{tá}\textit{tes} \textbf{e}jus.

    \item Conféssio et magnificéntia opus \textbf{e}jus:~* et justítia ejus manet in sǽ\textit{cu}\textit{lum} \textbf{sǽ}culi.
    
    \item Memóriam fecit mirabílium suórum,~† miséricors et miserátor \textbf{Dó}minus:~* escam dedit ti\textit{mén}\textit{ti}\textbf{bus} se.
    
    \item Memor erit in sǽculum testaménti \textbf{su}i:~* virtútem óperum suórum annuntiábit pó\textit{pu}\textit{lo} \textbf{su}o:
    
    \item Ut det illis hereditátem \textbf{gén}tium:~* ópera mánuum ejus véritas, \textit{et} \textit{ju}\textbf{dí}cium.
    
    \item Fidélia ómnia mandáta ejus:~† confirmáta in sǽculum \textbf{sǽ}culi,~* facta in veritáte et \textit{æ}\textit{qui}\textbf{tá}te.
    
    \item Redemptiónem misit pópulo \textbf{su}o:~* mandávit in ætérnum testa\textit{mén}\textit{tum} \textbf{su}um.
    
    \item Sanctum, et terríbile nomen \textbf{e}jus:~* inítium sapiéntiæ \textit{ti}\textit{mor} \textbf{Dó}mini.
    
    \item Intelléctus bonus ómnibus faciéntibus \textbf{e}um:~* laudátio ejus manet in sǽ\textit{cu}\textit{lum} \textbf{sǽ}culi.
    
    \item Glória Patri, et \textbf{Fí}lio,~* et Spirí\textit{tu}\textit{i} \textbf{Sanc}to.
    
    \item Sicut erat in princípio, et nunc, et \textbf{sem}per,~* et in sǽcula sæcu\textit{ló}\textit{rum}. \textbf{A}men.
  \end{enumerate}

  \medskip
  \grecommentary{\textit{Reprise de l'Antienne.}}
  \gabcsnippet{(c4) Ju(f)cun(h)dá(h!iwj!kvJ_I'_jwkjk)re(k_[oh:h+1mm]j_[oh:h+1mm]) (,) fí(ji)li(h)a(ij) Si(hg)on,(g.) (;) ex(g)súl(g')ta(g) sa(f_e)tis(d.) (,) fí(e')li(f)a(gh) Je(h')rú(g)sa(fg)lem,(gv_[oh:h]FD') al(f)le(fg)lú(g.){ia}.(g.) (::)}

  \newpage
  \vspace*{\fill}\
  \begin{normalsize}
    \begin{center}
      \par \textit{Grandes sont les œuvres du Seigneur ; tous ceux qui les aiment s'en instruisent.}
      \medskip
      \begin{enumerate}[label=\textcolor{red}{\emph{\arabic*}}]
        \item \textit{De tout cœur je rendrai grâce au Seigneur dans l'assemblée, parmi les justes.}
        \item \textit{Grandes sont les œuvres du Seigneur ; tous ceux qui les aiment s'en instruisent.}
        \item \textit{Noblesse et beauté dans ses actions : à jamais se maintiendra sa justice.}
        \item \textit{De ses merveilles il a laissé un mémorial ; le Seigneur est tendresse et pitié, il a donné des vivres à ses fidèles,}
        \item \textit{Gardant toujours mémoire de son
        alliance, il a montré sa force à son peuple.}
        \item \textit{Lui donnant le domaine des nations. Justesse et sûreté les œuvres de ses mains.}
        \item \textit{Sécurité, toutes ses lois, établies pour toujours et à jamais, accomplies avec droiture et sûreté ! }
        \item \textit{Il apporte la délivrance à son peuple ; son alliance est promulguée pour toujours.}
        \item \textit{Saint, redoutable est son nom, la sagesse commence avec la crainte du Seigneur.}
        \item \textit{Qui accomplit sa volonté en est éclairé. A jamais se maintiendra sa louange.}
        \item \textit{Gloire au Père, au Fils, et au Saint Esprit, }
        \item \textit{Comme il était au commencement, maintenant et toujours, et dans les siècles des siècles. Amen. }
      \end{enumerate}
    \end{center}
  \end{normalsize}
  \vspace*{\fill}\
  \newpage

  % ===== DEBUT Antienne =========
  \gresetinitiallines{1}
  \greillumination{\initfamily\fontsize{11mm}{11mm}\selectfont E}
  \gregorioscore{antiennes/an--ecce_dominus_veniet--solesmes_1961}
  \begin{center}
    \footnotesize{
      \textit{
        Voici que le Seigneur vient, et tous ses saints avec lui ; et il paraîtra une grande lumière en ce jour-là, alléluia. 
      }
    }
  \end{center}
  % \medskip
  % ===== FIN Antienne ===========

  % ===== DEBUT psaume ===========
  % gresetinitiallines : avec le parametre à 0, supprime l'ornement

  \gresetinitiallines{0}

  \begin{center}
    \normalsize{Psaume 111.}\\
  \end{center}
  % \smallskip
  \grechangedim{baselineskip}{50pt}{scalable}
  \gregorioscore{psaumes/psaume111-Va}
  \begin{enumerate}[label=\textcolor{red}{\emph{\arabic*}}]
    \setcounter{enumi}{1}
    \item Potens in terra erit semen \textbf{e}jus:~* generátio rectórum be\textbf{ne}di\textbf{cé}tur.

    \item Glória, et divítiæ in domo \textbf{e}jus:~* et justítia ejus manet in \textbf{sǽ}culum \textbf{sǽ}culi.
    
    \item Exórtum est in ténebris lumen \textbf{rec}tis:~* miséricors, et mise\textbf{rá}tor, et \textbf{jus}tus.
    
    \item Jucúndus homo qui miserétur et cómmodat,~† dispónet sermónes suos in ju\textbf{dí}cio:~* quia in ætérnum non \textbf{com}mo\textbf{vé}bitur.
    
    \item In memória ætérna erit \textbf{jus}tus:~* ab auditióne mala \textbf{non} ti\textbf{mé}bit.
    
    \item Parátum cor ejus speráre in Dómino,~† confirmátum est cor \textbf{e}jus:~* non commovébitur donec despíciat ini\textbf{mí}cos \textbf{su}os.
    
    \item Dispérsit, dedit paupéribus:~† justítia ejus manet in sǽculum \textbf{sǽ}culi,~* cornu ejus exaltábi\textbf{tur} in \textbf{gló}ria.
    
    \item Peccátor vidébit, et irascétur,~† déntibus suis fremet et ta\textbf{bé}scet:~* desidérium pecca\textbf{tó}rum per\textbf{í}bit.
    
    \item Glória Patri, et \textbf{Fí}lio,~* et Spi\textbf{rí}tui \textbf{Sanc}to.
    
    \item Sicut erat in princípio, et nunc, et \textbf{sem}per,~* et in sǽcula sæcu\textbf{ló}rum. \textbf{A}men.
  \end{enumerate}

  \medskip
  \grecommentary{\textit{Reprise de l'Antienne.}}
  \gabcsnippet{(c3) Ec(h_)ce(f'_) Dó(h)mi(g)nus(e_[uh:l]g) vé(f)ni(e)et,(d.) (,) et(d) o(e)mnes(d') san(e)cti(f') e(h)jus(i') cum(i) e(h)o :(h.) (;) et(h) e(h)rit(h) in(h') di(i)e(g') il(h)la(f.) (,) lux(d_e) ma(e)gna,(e') al(f)le(ef)lú(d.){ia}.(d.) (::)}

  \newpage
  \vspace*{\fill}\
  \begin{normalsize}
    \begin{center}
      \par \textit{Celui qui craint le Seigneur a une volonté ardente d’accomplir ses commandements.}
      \medskip
      \begin{enumerate}[label=\textcolor{red}{\emph{\arabic*}}]
        \item \textit{Heureux l’homme qui craint le Seigneur, qui aime entièrement sa volonté !}
        \item \textit{Sa lignée sera puissante sur la terre ; la
        race des justes est bénie.}
        \item \textit{Les richesses affluent dans sa maison : à
        jamais se maintiendra sa justice.}
        \item \textit{Lumière des cœurs droits, il s'est levé
        dans les ténèbres, l’homme de justice, de
        tendresse et de pitié.}
        \item \textit{L'homme de bien a pitié, il partage ; il
        mène ses affaires avec droiture, cet
        homme jamais ne tombera ;}
        \item \textit{Toujours on fera mémoire du juste, il ne
        craint pas l'annonce d'un malheur :}
        \item \textit{Le cœur ferme, il s'appuie sur le
        Seigneur. Son cœur est confiant, il ne
        craint pas : il verra ce que valaient ses
        oppresseurs.}
        \item \textit{A pleines mains, il donne au pauvre ; à
        jamais se maintiendra sa justice, sa
        puissance grandira, et sa gloire !}
        \item \textit{L'impie le voit et s'irrite ; il grince des
        dents et se détruit. L'ambition des impies
        se perdra.}
        \item \textit{Gloire au Père, au Fils, et au Saint Esprit, }
        \item \textit{Comme il était au commencement, maintenant et toujours, et dans les siècles des siècles. Amen. }
      \end{enumerate}
    \end{center}
  \end{normalsize}
  \vspace*{\fill}\
  \newpage

  % ===== DEBUT Antienne =========
  \gresetinitiallines{1}
  \greillumination{\initfamily\fontsize{11mm}{11mm}\selectfont O}
  \gregorioscore{antiennes/an--omnes_sitientes--solesmes_1961.1}
  \begin{center}
    \footnotesize{
      \textit{
        Vous tous qui avez soif, venez aux eaux, ; cherchez le Seigneur pendant qu’on peut le trouver, alléluia.
      }
    }
  \end{center}
  % ===== FIN Antienne ===========

  % ===== DEBUT psaume ===========
  % gresetinitiallines : avec le parametre à 0, supprime l'ornement
  \gresetinitiallines{0}

  \begin{center}
    \normalsize{Psaume 112.}\\
  \end{center}
  % \smallskip

  \gregorioscore{psaumes/psaume112-VIIC}

  \begin{enumerate}[label=\textcolor{red}{\emph{\arabic*}}]
    \setcounter{enumi}{1}
    \item Sit nomen Dómini \textbf{be}ne\textbf{díc}tum,~* ex hoc nunc, et \textbf{us}que in \textbf{sǽ}culum.

    \item A solis ortu usque \textbf{ad} oc\textbf{cá}sum,~* laudábile \textbf{no}men \textbf{Dó}mini.
    
    \item Excélsus super omnes \textbf{gen}tes \textbf{Dó}minus,~* et super cælos \textbf{gló}ria \textbf{e}jus.
    
    \item Quis sicut Dóminus, Deus noster, qui in \textbf{al}tis \textbf{há}bitat,~* et humília réspicit in cælo \textbf{et} in \textbf{ter}ra?
    
    \item Súscitans a \textbf{ter}ra \textbf{ín}opem,~* et de stércore \textbf{é}rigens \textbf{páu}perem:
    
    \item Ut cóllocet eum \textbf{cum} prin\textbf{cí}pibus,~* cum princípibus \textbf{pó}puli \textbf{su}i.
    
    \item Qui habitáre facit stéri\textbf{lem} in \textbf{do}mo,~* matrem fili\textbf{ó}rum læ\textbf{tán}tem.
    
    \item Glória \textbf{Pa}tri, et \textbf{Fí}lio,~* et Spi\textbf{rí}tui \textbf{Sanc}to.
    
    \item Sicut erat in princípio, et \textbf{nunc}, et \textbf{sem}per,~* et in sǽcula sæcu\textbf{ló}rum. \textbf{A}men.
  \end{enumerate}

  \medskip
  \grecommentary{\textit{Reprise de l'Antienne.}}
  \gabcsnippet{(c3) O(ig/ij)mnes(i'_[oh:h]) () si(i)ti(g)én(iji)tes,(h.) (,) ve(f')ní(h_)te(g) ad(fe) a(f)quas :(e.) (:) quaé(i_[uh:l]j!kwl)ri(k)te(j_kJH) Dó(j_k)mi(ji)num,(i.) (;) dum(ii) in(e)ve(g)ní(i)ri(h) pot(f_h)est,(gv_[oh:h]FE) al(fg)le(f)lú(e.){ia}.(e.)}

  \newpage
  \vspace*{\fill}\
  \begin{normalsize}
    \begin{center}
      \par \textit{Que le nom du Seigneur soit béni dans tous les siècles.}
      \medskip
      \begin{enumerate}[label=\textcolor{red}{\emph{\arabic*}}]
        \item \textit{Louez, serviteurs du Seigneur, louez le nom du Seigneur !}
        \item \textit{Béni soit le nom du Seigneur, maintenant et
        pour les siècles des siècles !}
        \item \textit{Du levant au couchant du soleil, loué soit le
        nom du Seigneur !}
        \item \textit{Le Seigneur domine tous les peuples, sa gloire
        domine les cieux.}
        \item \textit{Qui est semblable au Seigneur notre Dieu ?
        Lui, il siège là-haut, mais il abaisse son regard vers le ciel et vers la terre.}
        \item \textit{De la poussière il relève le faible, il retire le
        pauvre de la cendre}
        \item \textit{Pour qu'il siège parmi les princes, parmi les
        princes de son peuple.}
        \item \textit{Il installe en sa maison la femme stérile,
        heureuse mère au milieu de ses fils.}
        \item \textit{Gloire au Père, au Fils, et au Saint Esprit, }
        \item \textit{Comme il était au commencement, maintenant et toujours, et dans les siècles des siècles. Amen. }
      \end{enumerate}
    \end{center}
  \end{normalsize}
  \vspace*{\fill}\
  \newpage

  % ===== DEBUT Antienne =========
  \gresetinitiallines{1}
  \greillumination{\initfamily\fontsize{11mm}{11mm}\selectfont E}
  \gregorioscore{antiennes/an--ecce_veniet_propheta--solesmes_1961}
  \begin{center}
    \footnotesize{
      \textit{
        Voici que vient un grand prophète ; et c’est lui qui renouvellera Jérusalem, alléluia.
      }
    }
  \end{center}
  % ===== DEBUT psaume ===========
  % gresetinitiallines : avec le parametre à 0, supprime l'ornement
  \begin{center}
    \normalsize{Psaume 113.}
  \end{center}

  % gresetinitiallines : avec le parametre à 0, supprime l'ornement
  \gresetinitiallines{0}
  \gregorioscore{psaumes/psaume113-IVE}

  \begin{enumerate}[label=\textcolor{red}{\arabic*}]
    \setcounter{enumi}{1}
    \item Facta est Judǽa sanctificá\textit{ti}\textit{o} \textbf{e}jus,~* Israël \textit{pot}\textit{és}\textit{tas} \textbf{e}jus.

    \item Mare vi\textit{dit}, \textit{et} \textbf{fu}git:~* Jordánis convér\textit{sus} \textit{est} \textit{re}\textbf{trór}sum.
    
    \item Montes exsultavérunt \textit{ut} \textit{a}\textbf{rí}etes,~* et colles sic\textit{ut} \textit{a}\textit{gni} \textbf{ó}\textbf{vi}um.
    
    \item Quid est tibi, mare, \textit{quod} \textit{fu}\textbf{gís}ti:~* et tu, Jordánis, quia convér\textit{sus} \textit{es} \textit{re}\textbf{trór}sum?
    
    \item Montes, exsultástis sic\textit{ut} \textit{a}\textbf{rí}etes,~* et colles, sic\textit{ut} \textit{a}\textit{gni} \textbf{ó}\textbf{vi}um.
    
    \item A fácie Dómini mo\textit{ta} \textit{est} \textbf{ter}ra,~* a fáci\textit{e} \textit{De}\textit{i} \textbf{Ja}cob.
    
    \item Qui convértit petram in sta\textit{gna} \textit{a}\textbf{quá}rum,~* et rupem in \textit{fon}\textit{tes} \textit{a}\textbf{quá}rum.
    
    \item Non nobis, Dómi\textit{ne}, \textit{non} \textbf{no}bis:~* sed nómini \textit{tu}\textit{o} \textit{da} \textbf{gló}\textbf{ri}am.
    
    \item Super misericórdia tua, et veri\textit{tá}\textit{te} \textbf{tu}a:~* nequándo dicant gentes: Ubi est \textit{De}\textit{us} \textit{e}\textbf{ó}rum?
    
    \item Deus autem nos\textit{ter} \textit{in} \textbf{cæ}lo:~* ómnia quæcúmque \textit{vó}\textit{lu}\textit{it}, \textbf{fe}cit.
    
    \item Simulácra géntium argén\textit{tum}, \textit{et} \textbf{au}rum,~* ópera \textit{má}\textit{nu}\textit{um} \textbf{hó}\textbf{mi}num.
    
    \item Os habent, et \textit{non} \textit{lo}\textbf{quén}tur:~* óculos habent, \textit{et} \textit{non} \textit{vi}\textbf{dé}bunt.
    
    \item Aures habent, \textit{et} \textit{non} \textbf{áu}dient:~* nares habent, et \textit{non} \textit{o}\textit{do}\textbf{rá}bunt.
    
    \item Manus habent, et non palpábunt:~† pedes habent, et non \textit{am}\textit{bu}\textbf{lá}bunt:~* non clamábunt in \textit{gút}\textit{tu}\textit{re} \textbf{su}o.
    
    \item Símiles illis fiant qui fá\textit{ci}\textit{unt} \textbf{e}a:~* et omnes qui con\textit{fí}\textit{dunt} \textit{in} \textbf{e}is.
    
    \item Domus Israël sperá\textit{vit} \textit{in} \textbf{Dó}mino:~* adjútor eórum et pro\textit{téc}\textit{tor} \textit{e}\textbf{ó}\textbf{rum} est,
    
    \item Domus Aaron sperá\textit{vit} \textit{in} \textbf{Dó}mino:~* adjútor eórum et pro\textit{téc}\textit{tor} \textit{e}\textbf{ó}\textbf{rum} est,
    
    \item Qui timent Dóminum, speravé\textit{runt} \textit{in} \textbf{Dó}mino:~* adjútor eórum et pro\textit{téc}\textit{tor} \textit{e}\textbf{ó}\textbf{rum} est.
    
    \item Dóminus memor \textit{fu}\textit{it} \textbf{nos}tri:~* et be\textit{ne}\textit{dí}\textit{xit} \textbf{no}bis:
    
    \item Benedíxit dó\textit{mu}\textit{i} \textbf{Is}raël:~* benedíxit \textit{dó}\textit{mu}\textit{i} \textbf{A}\textbf{a}ron.
    
    \item Benedíxit ómnibus, qui \textit{ti}\textit{ment} \textbf{Dó}minum,~* pusíl\textit{lis} \textit{cum} \textit{ma}\textbf{jó}\textbf{ri}bus.
    
    \item Adjíciat Dó\textit{mi}\textit{nus} \textbf{su}per vos:~* super vos, et super \textit{fí}\textit{li}\textit{os} \textbf{ves}tros.
  \end{enumerate}

  \newpage
  \vspace*{\fill}\
  \begin{normalsize}
    \begin{center}
      \par \textit{Notre Dieu est dans le ciel, tout ce qu’il a voulu, il l’a fait.}
      \medskip
      \begin{enumerate}[label=\textcolor{red}{\emph{\arabic*}}]
        \item \textit{Quand Israël sortit d'Égypte, et Jacob, de chez un peuple étranger,}
        \item \textit{Juda fut pour Dieu un sanctuaire, Israël devint son domaine.}
        \item \textit{La mer voit et s'enfuit, le Jourdain retourne en arrière.}
        \item \textit{Comme des béliers, bondissent les montagnes, et les collines, comme des agneaux.}
        \item \textit{Qu'as-tu, mer, à t'enfuir, Jourdain, à retourner en arrière ?}
        \item \textit{Montagnes, pourquoi bondir comme des béliers, collines, comme des agneaux ?}
        \item \textit{Tremble, terre, devant le Maître, devant la face du Dieu de Jacob,}
        \item \textit{Lui qui change le rocher en source et la pierre en fontaine d’eau vive.}
        \item \textit{Non pas à nous, Seigneur, non pas à nous, mais à ton nom donne la gloire.}
        \item \textit{Pour ton amour et ta vérité.}
        \item \textit{Pourquoi les païens diraient-ils : « Où donc est leur Dieu ? »}
        \item \textit{Notre Dieu, il est au ciel ; tout ce qu'il veut, il le fait.}
        \item \textit{Leurs idoles : or et argent, ouvrages de mains humaines.} 
        \item \textit{Elles ont une bouche et ne parlent pas, des yeux et ne voient pas,}
        \item \textit{Des oreilles et n'entendent pas, des narines et ne sentent pas.}
        \item \textit{Leurs mains ne peuvent toucher, leurs pieds ne peuvent marcher, pas un son ne sort de leur gosier !}
        \item \textit{Qu'ils deviennent comme elles, tous ceux qui les font, ceux qui mettent leur foi en elles.}
        \item \textit{Israël, mets ta foi dans le Seigneur : le secours, le bouclier, c'est lui !}
        \item \textit{Famille d'Aaron, mets ta foi dans le Seigneur : le secours, le bouclier, c'est lui !}
        \item \textit{Vous qui le craignez, ayez foi dans le
        Seigneur : le secours, le bouclier, c'est lui !}
        \item \textit{Le Seigneur se souvient de nous : il bénira !
        Il bénira la famille d'Israël, et la famille
        d'Aaron}
        \item \textit{Il bénira tous ceux qui craignent le Seigneur,
        du plus grand au plus petit.}
        \item \textit{Que le Seigneur multiplie ses bienfaits pour
        vous et vos enfants !}
        \item \textit{Soyez bénis par le Seigneur qui a fait le ciel
        et la terre !}
        \item \textit{Le ciel, c'est le ciel du Seigneur ; aux
        hommes, il a donné la terre.}
        \item \textit{Les morts ne louent pas le Seigneur, ni ceux
        qui descendent au silence.}
        \item \textit{Nous, les vivants, bénissons le Seigneur,
        maintenant et pour les siècles des siècles}
        \item \textit{Gloire au Père, au Fils, et au Saint Esprit, }
        \item \textit{Comme il était au commencement, maintenant et toujours, et dans les siècles des siècles. Amen. }
      \end{enumerate}
    \end{center}
  \end{normalsize}
  \vspace*{\fill}\
  \newpage

  \begin{enumerate}[label=\textcolor{red}{\arabic*}]
    \setcounter{enumi}{23}
    \item Benedícti \textit{vos} \textit{a} \textbf{Dó}mino,~* qui fecit \textit{cæ}\textit{lum}, \textit{et} \textbf{ter}ram.

    \item Cælum \textit{cæ}\textit{li} \textbf{Dó}mino:~* terram autem dedit \textit{fí}\textit{li}\textit{is} \textbf{hó}\textbf{mi}num.
    
    \item Non mórtui laudá\textit{bunt} \textit{te}, \textbf{Dó}mine:~* neque omnes, qui descén\textit{dunt} \textit{in} \textit{in}\textbf{fér}num.
    
    \item Sed nos qui vívimus, benedí\textit{ci}\textit{mus} \textbf{Dó}mino,~* ex hoc nunc et \textit{us}\textit{que} \textit{in} \textbf{sǽ}\textbf{cu}lum.
    
    \item Glória Pa\textit{tri}, \textit{et} \textbf{Fí}lio,~* et Spi\textit{rí}\textit{tu}\textit{i} \textbf{Sanc}to.
    
    \item Sicut erat in princípio, et \textit{nunc}, \textit{et} \textbf{sem}per,~* et in sǽcula sæ\textit{cu}\textit{ló}\textit{rum}. \textbf{A}men.
  \end{enumerate}
  %  Répetition de l'Antienne
  \grecommentary{\textit{Reprise de l'Antienne.}}
  \gabcsnippet{(c3) Ec(e')ce(f) vé(h)ni(hi)et(i.) (,) Pro(h_i)phé(j)ta(ih) ma(i_[oh:h])gnus,(i.) (;) et(i) i(f')pse(i) re(g')no(h)vá(f_e)bit(f') (,) Je(e)rú(d')sa(e)lem,(gxf_g) al(h')le(h)lú(f.){ia}.(f.) (::)}

  \medskip
  \begin{center}
    \rule{2cm}{0.4pt}
  \end{center}
  \medskip

  \par \textit{\footnotesize\textcolor{red}{On se lève pour le Capitule.}}

  \begin{center}
    \textcolor{red}{\large{Capitule}}\\
    \small\textit{
      Épître aux Romains. 13, 11
    }
  \end{center}

  \begin{multicols}{2}
    \parindent=0pt
    Fratres : Hora est jam nos de somno súrgere. \textcolor{red}{†} Nunc enim própior est nostra salus, \textcolor{red}{*} quam cum credídimus. \\
    \textcolor{red}{\Rbar.} Deo grátias.

    \columnbreak

    \textit{ Frères : l’heure est déjà venue de nous réveiller de notre assoupissement, puisque nous sommes plus proches de notre salut que lorsque nous avons reçu la foi. \\
    \textcolor{red}{\Rbar.} Rendons grâce à Dieu.
    }
  \end{multicols}

  % \vspace*{\fill}\
  % \begin{center}
  %   \greseparator{3}{30}
  % \end{center}
  % \vspace*{\fill}\

  % \newpage

  \par \textit{\footnotesize\textcolor{red}{Le Célébrant entonne, ensuite, les Chantres et le Chœur alternent les versets. La Doxologie est chantée par tous.}}

  \begin{center}
    \textcolor{red}{\large{Hymne}}\\
  \end{center}
  
  \gresetinitiallines{1}
  \greillumination{\initfamily\fontsize{11mm}{11mm}\selectfont C}
  \gregorioscore{hymnes/hy--creator_alme_siderum--solesmes_1961}
  \bigskip
  % \begin{multicols}{2}
    \begin{normalsize}
      \begin{enumerate}[label=\textcolor{red}{\emph{\arabic*}}]
        \item \textit{Ô très doux créateur des cieux, Lampe éternelle des croyants, O Christ Sauveur de l’univers, Entends la voix des suppliants. }
        \item \textit{Pour ne point nous laisser périr par les ruses de Satan, cédantà votre amour, vous vous êters fait le remède de notre infirmité. }
        \item \textit{Pour expier l'universelle déchéance du monde, Victime immaculée, du sein d'une Vierge vous vous avancez vers la croix.}
        \item \textit{Oh vous qui, devant la puissance de votre gloire, à votre seul nom, voyez, au ciel et aux enfers, tout genou se plier tremblant.}
        \item \textit{Nous vous en prions, grand juge du dernier jour : par les armes de la céleste grâce, délivrez-nous de nos ennemis.}
        \item \textit{Louange, honneur, puissance et gloire A Dieu le Père et à son Fils ; A l’Esprit notre Défenseur, Dans les éternités sans fin. Amen.}
      \end{enumerate}
    \end{normalsize}
  % \end{multicols}

  \begin{center}
    \begin{footnotesize}
      \textcolor{red}{\textit{On chante le verset debout.}}
    \end{footnotesize}
    \begin{minipage}{0.8\linewidth}
      \gresetinitiallines{0}
      \gabcsnippet{(c3)<c><v>\Vbar</v>.</c> Ro(h)rá(h)te(h) cœ(h)li(h) dé(h)su(h)per(h), et(h) nu(h)bes(h) plu(h)ant(h) jus(h)tum. (g'_) (hvGF'Efgf.) (::) (Z) <c><v>\Rbar</v>.</c> Ape(h)riá(h)tur(h) te(h)rra(h) et(h) gér(h)mi(h)net(h) Sal(h)va(h)tó(h)rem. (g'_) (hvGF'Efgf.) (::)}
      \bigskip
      \begin{center}
        \textit{\textcolor{red}{\Vbar.} Cieux, répandez votre rosée, et vous, nuées, faites pleuvoir le juste.}\\
        \textit{\textcolor{red}{\Rbar.} Que la terre s’ouvre et qu’elle germe le Sauveur.}
      \end{center}
    \end{minipage}
  \end{center}
  \normalsize

  \vspace*{\fill}\
  \begin{center}
    \greseparator{3}{30}
  \end{center}
  \vspace*{\fill}\

  \newpage

  \begin{center}
    \textcolor{red}{\large{Antienne à Magnificat}}\\
  \end{center}

  \gresetinitiallines{1}
  \greillumination{\initfamily\fontsize{11mm}{11mm}\selectfont N}
  \gregorioscore{antiennes/an--ne_timeas_maria--solesmes}
  \medskip
  \begin{center}
    \footnotesize{\textit{
      Ne craignez point, Marie, vous avez trouvé grâce auprès du Seigneur ; voici que vous concevrez, et enfanterez un fils, alléluia. 
    }}
  \end{center}
  \medskip

  \gresetinitiallines{0}
  \gregorioscore{magnificat/magnificat-VIIIG}

  \begin{enumerate}[label=\textcolor{red}{\arabic*}]
    \setcounter{enumi}{2}
    \item Quia respéxit humilitátem ancíllæ \textbf{su}æ:~* ecce enim ex hoc beátam me dicent omnes gene\textit{ra}\textit{ti}\textbf{ó}nes.

    \item Quia fecit mihi magna qui \textbf{pot}ens est:~* et sanctum \textit{no}\textit{men} \textbf{e}jus.
    
    \item Et misericórdia ejus a progénie in pro\textbf{gé}nies~* timén\textit{ti}\textit{bus} \textbf{e}um.
    
    \item Fecit poténtiam in bráchio \textbf{su}o:~* dispérsit supérbos mente \textit{cor}\textit{dis} \textbf{su}i.
    
    \item Depósuit poténtes de \textbf{se}de,~* et exal\textit{tá}\textit{vit} \textbf{hú}miles.
    
    \item Esuriéntes implévit \textbf{bo}nis:~* et dívites dimí\textit{sit} \textit{in}\textbf{á}nes.
    
    \item Suscépit Israël púerum \textbf{su}um,~* recordátus misericór\textit{di}\textit{æ} \textbf{su}æ.
    
    \item Sicut locútus est ad patres \textbf{nos}tros,~* Abraham et sémini e\textit{jus} \textit{in} \textbf{sǽ}cula.
    \begin{center}
      \textit{\footnotesize \textcolor{red}{On attend la fin de l'encensement avant de chanter la doxologie.}}
    \end{center}
    \item Glória Patri, et \textbf{Fí}lio,~* et Spirí\textit{tu}\textit{i} \textbf{Sanc}to.
    
    \item Sicut erat in princípio, et nunc, et \textbf{sem}per,~* et in sǽcula sæcu\textit{ló}\textit{rum}. \textbf{A}men.
  \end{enumerate}

  % \grecommentary{\textit{Reprise de l'Antienne.}}
  \gabcsnippet{(c4) Ne(g) tí(gj)me(hg)as,(h') () Ma(f)rí(g_h)a,(g_h) (,) in(g)ve(g)ní(h')sti(g) grá(h_j!kv)ti(ji)am(g_h) a(i)pud(h') Dó(g)mi(fg)num:(g.) (:) ec(h')ce(g) con(ge~)cí(f!g'h)pi(fe)es,(d.) (;) et(d) pá(e_f)ri(gh)es(h) fí(g')li(g)um,(gv_FD') al(f)le(fg)lú(g.)ia.(g.) (::)}

  \newpage

  \begin{center}
    \textcolor{red}{\large{Oraison}}\\
  \end{center}

  \begin{multicols}{2}
    \parindent=0pt
    \begin{flushright}
      \textcolor{red}{\Vbar.} Dominus vobiscum.\\
      \textcolor{red}{\Rbar.} Et cum spiritu tuo.\\
    \end{flushright}

    \columnbreak
    
    \textit{\textcolor{red}{\Vbar.} Le Seigneur soit avec vous.\\
    \textcolor{red}{\Rbar.} Et avec votre esprit.}\\
  \end{multicols}

  \begin{multicols}{2}
    \parindent=0pt
    Excita, quæsumus, Dómine, poténtiam tuam, et veni :  \textcolor{red}{†} ut ab imminéntibus peccatórum nostrórum perículis, te mereámur protegénte éripi,  \textcolor{red}{*} te liberánte salvári : \\ Qui vivis et regnas cum Deo Patre in unitáte Spíritus Sancti Deus : per ómnia sæcula sæculórum. 
    \textcolor{red}{\Rbar.} Amen.

    \columnbreak

    \textit{ Réveillez, Seigneur, votre puissance, et venez : pour que nous méritions d’être arrachés aux périls imminents auxquels nous exposent nos péchés, et que nous soyons sauvés par votre action libératrice. Vous qui vivez et règnez avec Dieu le Père en l’unité du Saint Esprit, Dieu pour tous les siècles des siècles. 
    Amen.
    }
  \end{multicols}

  \medskip

  \begin{center}
    \textcolor{red}{\large{Conclusion de l'office}}
  \end{center}
  
  
  \begin{multicols}{2}
    \parindent=0pt
    \begin{flushright}
      \textcolor{red}{\Vbar.} Dominus vobiscum.\\
      \textcolor{red}{\Rbar.} Et cum spiritu tuo.\\
    \end{flushright}
  
    \columnbreak
    
    \textit{\textcolor{red}{\Vbar.} Le Seigneur soit avec vous.\\
    \textcolor{red}{\Rbar.} Et avec votre esprit.}\\
  \end{multicols}
  \bigskip
  \gresetinitiallines{1}
  \greillumination{\initfamily\fontsize{11mm}{11mm}\selectfont B}
  \gregorioscore{or--benedicamus_domino_(sundays_of_advent_and_lent)--solesmes_1961}
  \begin{center}
    \begin{footnotesize}
      \textcolor{red}{\textit{Sur un ton très grave : }}
    \end{footnotesize}
  \end{center}
  \begin{multicols}{2}
    \parindent=0pt
    \textcolor{red}{\Vbar.} Fidélium ánimæ per misericórdiam Dei requiéscant in pace.\\
    \textcolor{red}{\Rbar.} Amen.\\

    \columnbreak
    
    \textit{\textcolor{red}{\Vbar.} Que les âmes des fidèles défunts, par la
    miséricorde de Dieu, reposent en paix.\\
    \textcolor{red}{\Rbar.} Amen.}\\
  \end{multicols}

  \vspace*{\fill}\
  \begin{center}
    \greseparator{3}{30}
  \end{center}
  \vspace*{\fill}\

  \newpage

  \begin{center}
    \begin{LARGE}
      II\textsuperscript{e} dimanche de l'Avent
    \end{LARGE}
  \end{center}
  \medskip

  % ===== DEBUT Antienne =========
  \greillumination{\initfamily\fontsize{11mm}{11mm}\selectfont E}
  \gregorioscore{antiennes/an--ecce_in_nubibus--solesmes_1961}
  \begin{center}
    \footnotesize{
      \textit{
        Voici que dans les nuées du ciel le Seigneur viendra avec une grande puissance, alléluia. 
      }
    }
  \end{center}
  % \medskip
  % ===== FIN Antienne ===========

  % ===== DEBUT psaume ===========
  % gresetinitiallines : avec le parametre à 0, supprime l'ornement
  \gresetinitiallines{0}

  \begin{center}
    \normalsize{Psaume 109.}
  \end{center}

  \gregorioscore{psaumes/psaume109-Ig}
  \begin{enumerate}[label=\textcolor{red}{\emph{\arabic*}}]
    \setcounter{enumi}{2}
    \item Donec ponam ini\textbf{mí}cos \textbf{tu}os,~* scabéllum pe\textit{dum} \textit{tu}\textbf{ó}rum.

    \item Virgam virtútis tuæ emíttet Dómi\textbf{nus} ex \textbf{Si}on:~* domináre in médio inimicó\textit{rum} \textit{tu}\textbf{ó}rum.
    
    \item Tecum princípium in die virtútis tuæ in splendóri\textbf{bus} sanc\textbf{tó}rum:~* ex útero ante lucíferum \textit{gé}\textit{nu}\textbf{i} te.
    
    \item Jurávit Dóminus, et non pœni\textbf{té}bit \textbf{e}um:~* Tu es sacérdos in ætérnum secúndum órdi\textit{nem} \textit{Mel}\textbf{chí}sedech.
    
    \item Dóminus a \textbf{dex}tris \textbf{tu}is,~* confrégit in die iræ \textit{su}\textit{æ} \textbf{re}ges.
    
    \item Judicábit in natiónibus, im\textbf{plé}bit ru\textbf{í}nas:~* conquassábit cápita in ter\textit{ra} \textit{mul}\textbf{tó}rum.
    
    \item De torrénte in \textbf{vi}a \textbf{bi}bet:~* proptérea exal\textit{tá}\textit{bit} \textbf{ca}put.
    
    \item Glória \textbf{Pa}tri, et \textbf{Fí}lio,~* et Spirí\textit{tu}\textit{i} \textbf{Sanc}to.
    
    \item Sicut erat in princípio, et \textbf{nunc}, et \textbf{sem}per,~* et in sǽcula sæcu\textit{ló}\textit{rum}. \textbf{A}men.
  \end{enumerate}
  \medskip
  \grecommentary{\textit{Reprise de l'Antienne.}}
  \gabcsnippet{(c4) Ec(dd)ce(c) in(d) nú(e_[uh:l]f)bi(g)bus(fe) cae(d)li(d_c) (,) <c>*</c> Dó(f_g)mi(h)nus(g') vé(h)ni(gf)et,(f.) (;) cum(f) pot(e')e(f)stá(g')te(f) ma(fe~)gna,(dc) al(d)le(fe)lú(d.){ia}.(d.) (::)}

  \newpage

  % ===== DEBUT Antienne =========
  \gresetinitiallines{1}
  \greillumination{\initfamily\fontsize{11mm}{11mm}\selectfont U}
  \gregorioscore{antiennes/an--urbs_fortitudinis--solesmes_1961.1}
  \begin{center}
    \footnotesize{
      \textit{
        Notre ville forte est Sion ; le Sauveur y sera mis comme mur et avant-mur. Ouvrez les portes, parce que Dieu est avec nous, alléluia.
      }
    }
  \end{center}
  % \medskip
  % ===== FIN Antienne ===========

  % ===== DEBUT psaume ===========
  % gresetinitiallines : avec le parametre à 0, supprime l'ornement
  \gresetinitiallines{0}

  \begin{center}
    \normalsize{Psaume 110.}
  \end{center}

  \gregorioscore{psaumes/psaume110-VIId}

  \begin{enumerate}[label=\textcolor{red}{\emph{\arabic*}}]
    \setcounter{enumi}{1}
    \item Magna \textbf{ó}pera \textbf{Dó}mini:~* exquisíta in omnes volun\textbf{tá}tes \textbf{e}jus.

    \item Conféssio et magnificéntia \textbf{o}pus \textbf{e}jus:~* et justítia ejus manet in \textbf{sǽ}culum \textbf{sǽ}culi.
    
    \item Memóriam fecit mirabílium suórum,~† miséricors et mise\textbf{rá}tor \textbf{Dó}minus:~* escam dedit ti\textbf{mén}ti\textbf{bus} se.
    
    \item Memor erit in sǽculum testa\textbf{mén}ti \textbf{su}i:~* virtútem óperum suórum annuntiábit \textbf{pó}pulo \textbf{su}o:
    
    \item Ut det illis heredi\textbf{tá}tem \textbf{gén}tium:~* ópera mánuum ejus véritas, \textbf{et} ju\textbf{dí}cium.
    
    \item Fidélia ómnia mandáta ejus:~† confirmáta in \textbf{sǽ}culum \textbf{sǽ}culi,~* facta in veritáte et \textbf{æ}qui\textbf{tá}te.
    
    \item Redemptiónem misit \textbf{pó}pulo \textbf{su}o:~* mandávit in ætérnum testa\textbf{mén}tum \textbf{su}um.
    
    \item Sanctum, et terríbile \textbf{no}men \textbf{e}jus:~* inítium sapiéntiæ \textbf{ti}mor \textbf{Dó}mini.
    
    \item Intelléctus bonus ómnibus faci\textbf{én}tibus \textbf{e}um:~* laudátio ejus manet in \textbf{sǽ}culum \textbf{sǽ}culi.
    
    \item Glória \textbf{Pa}tri, et \textbf{Fí}lio,~* et Spi\textbf{rí}tui \textbf{Sanc}to.
    
    \item Sicut erat in princípio, et \textbf{nunc}, et \textbf{sem}per,~* et in sǽcula sæcu\textbf{ló}rum. \textbf{A}men.
  \end{enumerate}

  \medskip
  \grecommentary{\textit{Reprise de l'Antienne.}}
  \gabcsnippet{(c2) U{r}bs(c_g) for(g)ti(g)tú(g')di(g)nis(g') no(i)strae(g) Si(fg)on,(g.) (;) Sal(exe')vá(g)tor(f') po(exg)né(fe)tur(d') in(exe) e(f)a(exe[ll:1]d) mu(c)rus(c.) (;) et(ghg___) an(gf~)te(g)mu(gj)rá(jjg)le :(g.) (:) a(g)pe(g')rí(h)te(gf) por(g)tas,(e.[ll:1]d!ewf.) (;) qui(f)a(ffe) no(d)bís(e_[uh:l]f)cum(d) De(dc)us,(b) al(cd)le(d)lú(c.){ia}.(c.) (::)}

  \newpage

  % ===== DEBUT Antienne =========
  \gresetinitiallines{1}
  \greillumination{\initfamily\fontsize{11mm}{11mm}\selectfont E}
  \gregorioscore{antiennes/an--ecce_apparebit--solesmes_1961}
  \begin{center}
    \footnotesize{
      \textit{
        Voici que paraîtra le Seigneur et il ne trompera pas ; s’il met un délai, attends le, car il viendra et il ne tardera pas, alléluia.
      }
    }
  \end{center}
  % \medskip
  % ===== FIN Antienne ===========

  % ===== DEBUT psaume ===========
  % gresetinitiallines : avec le parametre à 0, supprime l'ornement

  \gresetinitiallines{0}

  \begin{center}
    \normalsize{Psaume 111.}\\
  \end{center}
  % \smallskip
  \grechangedim{baselineskip}{50pt}{scalable}
  \gregorioscore{psaumes/psaume111-VIIa}
  \begin{enumerate}[label=\textcolor{red}{\emph{\arabic*}}]
    \setcounter{enumi}{1}
    \item Potens in terra erit \textbf{se}men \textbf{e}jus:~* generátio rectórum be\textbf{ne}di\textbf{cé}tur.

    \item Glória, et divítiæ in \textbf{do}mo \textbf{e}jus:~* et justítia ejus manet in \textbf{sǽ}culum \textbf{sǽ}culi.
    
    \item Exórtum est in ténebris \textbf{lu}men \textbf{rec}tis:~* miséricors, et mise\textbf{rá}tor, et \textbf{jus}tus.
    
    \item Jucúndus homo qui miserétur et cómmodat,~† dispónet sermónes suos \textbf{in} ju\textbf{dí}cio:~* quia in ætérnum non \textbf{com}mo\textbf{vé}bitur.
    
    \item In memória ætérna \textbf{e}rit \textbf{jus}tus:~* ab auditióne mala \textbf{non} ti\textbf{mé}bit.
    
    \item Parátum cor ejus speráre in Dómino,~† confirmátum \textbf{est} cor \textbf{e}jus:~* non commovébitur donec despíciat ini\textbf{mí}cos \textbf{su}os.
    
    \item Dispérsit, dedit paupéribus:~† justítia ejus manet in \textbf{sǽ}culum \textbf{sǽ}culi,~* cornu ejus exaltábi\textbf{tur} in \textbf{gló}ria.
    
    \item Peccátor vidébit, et irascétur,~† déntibus suis fremet \textbf{et} ta\textbf{bé}scet:~* desidérium pecca\textbf{tó}rum per\textbf{í}bit.
    
    \item Glória \textbf{Pa}tri, et \textbf{Fí}lio,~* et Spi\textbf{rí}tui \textbf{Sanc}to.
    
    \item Sicut erat in princípio, et \textbf{nunc}, et \textbf{sem}per,~* et in sǽcula sæcu\textbf{ló}rum. \textbf{A}men.
  \end{enumerate}

  \medskip
  \grecommentary{\textit{Reprise de l'Antienne.}}
  \gabcsnippet{(c3) Ec(e')ce(e) ap(g')pa(h)ré(ih/ij)bit(i) Dó(i_[uh:l]j)mi(ih)nus,(h'_) (,) et(h) non(h) men(h)ti(h)é(gf)tur :(ef..) (:) si(d') mo(e)ram(f') fé(h)ce(hi)rit,(i.) (,) ex(g)spé(h)cta(f') e(f)um,(e.) (;) qui(e)a(e) vé(e_[uh:l]f)ni(dc)et,(b'_) (,) et(d) non(ef) tar(f')dá(e)bit,(ev_[oh:h]DB') al(d)le(de)lú(e.){ia}.(e.) (::)}

  \newpage

  % ===== DEBUT Antienne =========
  \gresetinitiallines{1}
  \greillumination{\initfamily\fontsize{11mm}{11mm}\selectfont M}
  \gregorioscore{antiennes/an--montes_et_colles--solesmes_1961}
  \begin{center}
    \footnotesize{
      \textit{
        Les montagnes et les collines chanteront devant Dieu des louanges, et tous les arbres des forêts battront des mains, parce que le Seigneur dominateur viendra pour régner éternellement, alléluia, alléluia. 
      }
    }
  \end{center}
  % ===== FIN Antienne ===========

  % ===== DEBUT psaume ===========
  % gresetinitiallines : avec le parametre à 0, supprime l'ornement
  \gresetinitiallines{0}

  \begin{center}
    \normalsize{Psaume 112.}\\
  \end{center}
  % \smallskip

  \gregorioscore{psaumes/psaume112-If}

  \begin{enumerate}[label=\textcolor{red}{\emph{\arabic*}}]
    \setcounter{enumi}{1}
    \item Sit nomen Dómini \textbf{be}ne\textbf{díc}tum,~* ex hoc nunc, et us\textit{que} \textit{in} \textbf{sǽ}culum.

    \item A solis ortu usque \textbf{ad} oc\textbf{cá}sum,~* laudábile \textit{no}\textit{men} \textbf{Dó}mini.
    
    \item Excélsus super omnes \textbf{gen}tes \textbf{Dó}minus,~* et super cælos gló\textit{ri}\textit{a} \textbf{e}jus.
    
    \item Quis sicut Dóminus, Deus noster, qui in \textbf{al}tis \textbf{há}bitat,~* et humília réspicit in cælo \textit{et} \textit{in} \textbf{ter}ra?
    
    \item Súscitans a \textbf{ter}ra \textbf{ín}opem,~* et de stércore é\textit{ri}\textit{gens} \textbf{páu}perem:
    
    \item Ut cóllocet eum \textbf{cum} prin\textbf{cí}pibus,~* cum princípibus pó\textit{pu}\textit{li} \textbf{su}i.
    
    \item Qui habitáre facit stéri\textbf{lem} in \textbf{do}mo,~* matrem filió\textit{rum} \textit{læ}\textbf{tán}tem.
    
    \item Glória \textbf{Pa}tri, et \textbf{Fí}lio,~* et Spirí\textit{tu}\textit{i} \textbf{Sanc}to.
    
    \item Sicut erat in princípio, et \textbf{nunc}, et \textbf{sem}per,~* et in sǽcula sæcu\textit{ló}\textit{rum}. \textbf{A}men.
  \end{enumerate}

  \medskip
  \grecommentary{\textit{Reprise de l'Antienne.}}
  \gabcsnippet{(c4) Mon(d_c~)tes(f) et(g) col(f_h)les(h.) (,) can(hg~)tá(h)bunt(g) co(e')ram(f) De(g)o(fe) lau(d)dem,(d.) (;) et(d) ó(f)mni(fg)a(g') li(h)gna(g) sil(fe~)vá(f_g)rum(g.) (,) plau(hj~)dent(ji~) má(g)ni(ij)bus :(h.) (:) quó(f)ni(gh)am(h_g) vé(e)ni(fg)et(g_[oh:h]f) Dó(d')mi(f)nus(e') Do(f)mi(d)ná(d)tor(c'_[oh:h]) (,) in(d) re(dd)gnum(c) ae(e)tér(ge~)num,(f_g) (,) al(e)le(c')lú(e){ia},(g') al(e)le(fe)lú(d.){ia}.(d.) (::) }

  \newpage

  % ===== DEBUT Antienne =========
  \gresetinitiallines{1}
  \greillumination{\initfamily\fontsize{11mm}{11mm}\selectfont E}
  \gregorioscore{antiennes/an--ecce_dominus_noster..._ut_illuminet--solesmes_1961}
  \begin{center}
    \footnotesize{
      \textit{
        Voici que notre Seigneur viendra avec puissance et illuminera les yeux de ses serviteurs, alléluia.
      }
    }
  \end{center}
  % ===== DEBUT psaume ===========
  % gresetinitiallines : avec le parametre à 0, supprime l'ornement
  \begin{center}
    \normalsize{Psaume 113.}
  \end{center}

  % gresetinitiallines : avec le parametre à 0, supprime l'ornement
  \gresetinitiallines{0}
  \gregorioscore{psaumes/psaume113-IIIa}

  \begin{enumerate}[label=\textcolor{red}{\arabic*}]
    \setcounter{enumi}{1}
    \item Facta est Judǽa sanctifi\textbf{cá}tio \textbf{e}jus,~* Israël potés\textit{tas} \textbf{e}jus.

    \item Mare \textbf{vi}dit, et \textbf{fu}git:~* Jordánis convérsus est \textit{re}\textbf{trór}sum.
    
    \item Montes exsultavérunt \textbf{ut} a\textbf{rí}\textbf{e}tes,~* et colles sicut a\textit{gni} \textbf{ó}vium.
    
    \item Quid est tibi, mare, \textbf{quod} fu\textbf{gís}ti:~* et tu, Jordánis, quia convérsus es \textit{re}\textbf{trór}sum?
    
    \item Montes, exsultástis \textbf{sic}ut a\textbf{rí}\textbf{e}tes,~* et colles, sicut a\textit{gni} \textbf{ó}vium.
    
    \item A fácie Dómini \textbf{mo}ta est \textbf{ter}ra,~* a fácie De\textit{i} \textbf{Ja}cob.
    
    \item Qui convértit petram in \textbf{sta}gna a\textbf{quá}rum,~* et rupem in fontes \textit{a}\textbf{quá}rum.
    
    \item Non nobis, Dómi\textbf{ne}, non \textbf{no}bis:~* sed nómini tuo \textit{da} \textbf{gló}riam.
    
    \item Super misericórdia tua, et veri\textbf{tá}te \textbf{tu}a:~* nequándo dicant gentes: Ubi est Deus \textit{e}\textbf{ó}rum?
    
    \item Deus autem \textbf{nos}ter in \textbf{cæ}lo:~* ómnia quæcúmque vólu\textit{it}, \textbf{fe}cit.
    
    \item Simulácra géntium ar\textbf{gén}tum, et \textbf{au}rum,~* ópera mánu\textit{um} \textbf{hó}minum.
    
    \item Os habent, et \textbf{non} lo\textbf{quén}tur:~* óculos habent, et non \textit{vi}\textbf{dé}bunt.
    
    \item Aures habent, \textbf{et} non \textbf{áu}\textbf{di}ent:~* nares habent, et non o\textit{do}\textbf{rá}bunt.
    
    \item Manus habent, et non palpábunt:~† pedes habent, et non \textbf{am}bu\textbf{lá}bunt:~* non clamábunt in gúttu\textit{re} \textbf{su}o.
    
    \item Símiles illis fiant qui \textbf{fá}ciunt \textbf{e}a:~* et omnes qui confídunt \textit{in} \textbf{e}is.
    
    \item Domus Israël spe\textbf{rá}vit in \textbf{Dó}\textbf{mi}no:~* adjútor eórum et protéctor \textit{e}\textbf{ó}rum est,
    
    \item Domus Aaron spe\textbf{rá}vit in \textbf{Dó}\textbf{mi}no:~* adjútor eórum et protéctor \textit{e}\textbf{ó}rum est,
    
    \item Qui timent Dóminum, spera\textbf{vé}runt in \textbf{Dó}\textbf{mi}no:~* adjútor eórum et protéctor \textit{e}\textbf{ó}rum est.
    
    \item Dóminus memor \textbf{fu}it \textbf{nos}tri:~* et benedí\textit{xit} \textbf{no}bis:
    
    \item Benedíxit \textbf{dó}mui \textbf{Is}\textbf{ra}ël:~* benedíxit dómu\textit{i} \textbf{A}aron.
    
    \item Benedíxit ómnibus, qui \textbf{ti}ment \textbf{Dó}\textbf{mi}num,~* pusíllis cum \textit{ma}\textbf{jó}ribus.
    
    \item Adjíciat \textbf{Dó}minus \textbf{su}\textbf{per} vos:~* super vos, et super fíli\textit{os} \textbf{ves}tros.
    
    \item Benedícti \textbf{vos} a \textbf{Dó}\textbf{mi}no,~* qui fecit cælum, \textit{et} \textbf{ter}ram.
    
    \item Cælum \textbf{cæ}li \textbf{Dó}\textbf{mi}no:~* terram autem dedit fíli\textit{is} \textbf{hó}minum.
    
    \item Non mórtui lau\textbf{dá}bunt te, \textbf{Dó}\textbf{mi}ne:~* neque omnes, qui descéndunt in \textit{in}\textbf{fér}num.
    
    \item Sed nos qui vívimus, bene\textbf{dí}cimus \textbf{Dó}\textbf{mi}no,~* ex hoc nunc et usque \textit{in} \textbf{sǽ}culum.
    
    \item Glória \textbf{Pa}tri, et \textbf{Fí}\textbf{li}o,~* et Spirítu\textit{i} \textbf{Sanc}to.
    
    \item Sicut erat in princípio, et \textbf{nunc}, et \textbf{sem}per,~* et in sǽcula sæculó\textit{rum}. \textbf{A}men.
  \end{enumerate}

  %  Répetition de l'Antienne
  \grecommentary{\textit{Reprise de l'Antienne.}}
  \gabcsnippet{(c4) Ec(g)ce(g) Dó(hi)mi(i)nus(i) no(i)ster(i'_[oh:h]) <c>*</c>(,) cum(i) vir(i')tú(j)te(h') vé(i)ni(hg)et,(g.) (;) ut(g) il(h)lú(j')mi(i)net(g') ó(h)cu(g)los(g'_[oh:h]) (,) ser(g)vó(g_[uh:l]h)rum(f') su(g)ó(f)rum,(f_d) al(e_[uh:l]f)le(g)lú(gg){ia}.(e.) (::)}

  \medskip
  \begin{center}
    \rule{2cm}{0.4pt}
  \end{center}
  \medskip

  \par \textit{\footnotesize\textcolor{red}{On se lève pour le Capitule.}}

  \begin{center}
    \textcolor{red}{\large{Capitule}}\\
    \small\textit{
      Épître aux Romains. 14, 4
    }
  \end{center}

  \begin{multicols}{2}
    \parindent=0pt
    Fratres : Quæcúmque scripta sunt, ad nostram doctrínam scripta sunt :  \textcolor{red}{†} ut per patiéntiam, et consolatiónem Scripturárum, \textcolor{red}{*} spem habeámus.  \\
    \textcolor{red}{\Rbar.} Deo grátias.

    \columnbreak

    \textit{ Frères : Car tout ce qui est écrit, a été écrit pour notre instruction ; afin que nous concevions l’espérance par la patience, et par la consolation que les Écritures nous donnent.  \\
    \textcolor{red}{\Rbar.} Rendons grâce à Dieu.
    }
  \end{multicols}


  \par \textit{\footnotesize\textcolor{red}{Le Célébrant entonne, ensuite, les Chantres et le Chœur alternent les versets. La Doxologie est chantée par tous.}}

  \begin{center}
    \textcolor{red}{\large{Hymne}}\\
  \end{center}
  
  \gresetinitiallines{1}
  \greillumination{\initfamily\fontsize{11mm}{11mm}\selectfont C}
  \gregorioscore{hymnes/hy--creator_alme_siderum--solesmes_1961}
  \bigskip

  \begin{center}
    \begin{footnotesize}
      \textcolor{red}{\textit{On chante le verset debout.}}
    \end{footnotesize}
    \begin{minipage}{0.8\linewidth}
      \gresetinitiallines{0}
      \gabcsnippet{(c3)<c><v>\Vbar</v>.</c> Ro(h)rá(h)te(h) cœ(h)li(h) dé(h)su(h)per(h), et(h) nu(h)bes(h) plu(h)ant(h) jus(h)tum. (g'_) (hvGF'Efgf.) (::) (Z) <c><v>\Rbar</v>.</c> Ape(h)riá(h)tur(h) te(h)rra(h) et(h) gér(h)mi(h)net(h) Sal(h)va(h)tó(h)rem. (g'_) (hvGF'Efgf.) (::)}
      \bigskip
      \begin{center}
        \textit{\textcolor{red}{\Vbar.} Cieux, répandez votre rosée, et vous, nuées, faites pleuvoir le juste.}\\
        \textit{\textcolor{red}{\Rbar.} Que la terre s’ouvre et qu’elle germe le Sauveur.}
      \end{center}
    \end{minipage}
  \end{center}
  \normalsize

  \vspace*{\fill}\
  \begin{center}
    \greseparator{3}{30}
  \end{center}
  \vspace*{\fill}\

  \newpage

  \begin{center}
    \textcolor{red}{\large{Antienne à Magnificat}}\\
  \end{center}

  \gresetinitiallines{1}
  \greillumination{\initfamily\fontsize{11mm}{11mm}\selectfont T}
  \gregorioscore{antiennes/an--tu_es_qui_venturus_es..._an_alium--solesmes_1961}
  \medskip
  \begin{center}
    \footnotesize{\textit{
      Ne craignez point, Marie, vous avez trouvé grâce auprès du Seigneur ; voici que vous concevrez, et enfanterez un fils, alléluia. 
    }}
  \end{center}
  \medskip

  \gresetinitiallines{0}
  \gregorioscore{magnificat/magnificat-VIIIG}

  \begin{enumerate}[label=\textcolor{red}{\arabic*}]
    \setcounter{enumi}{2}
    \item Quia respéxit humilitátem ancíllæ \textbf{su}æ:~* ecce enim ex hoc beátam me dicent omnes gene\textit{ra}\textit{ti}\textbf{ó}nes.

    \item Quia fecit mihi magna qui \textbf{pot}ens est:~* et sanctum \textit{no}\textit{men} \textbf{e}jus.
    
    \item Et misericórdia ejus a progénie in pro\textbf{gé}nies~* timén\textit{ti}\textit{bus} \textbf{e}um.
    
    \item Fecit poténtiam in bráchio \textbf{su}o:~* dispérsit supérbos mente \textit{cor}\textit{dis} \textbf{su}i.
    
    \item Depósuit poténtes de \textbf{se}de,~* et exal\textit{tá}\textit{vit} \textbf{hú}miles.
    
    \item Esuriéntes implévit \textbf{bo}nis:~* et dívites dimí\textit{sit} \textit{in}\textbf{á}nes.
    
    \item Suscépit Israël púerum \textbf{su}um,~* recordátus misericór\textit{di}\textit{æ} \textbf{su}æ.
    
    \item Sicut locútus est ad patres \textbf{nos}tros,~* Abraham et sémini e\textit{jus} \textit{in} \textbf{sǽ}cula.
    \begin{center}
      \textit{\footnotesize \textcolor{red}{On attend la fin de l'encensement avant de chanter la doxologie.}}
    \end{center}
    \item Glória Patri, et \textbf{Fí}lio,~* et Spirí\textit{tu}\textit{i} \textbf{Sanc}to.
    
    \item Sicut erat in princípio, et nunc, et \textbf{sem}per,~* et in sǽcula sæcu\textit{ló}\textit{rum}. \textbf{A}men.
  \end{enumerate}

  % \grecommentary{\textit{Reprise de l'Antienne.}}
  \gabcsnippet{(c4) Tu(f_g) es(g'_[oh:h]) qui(h) ven(g')tú(j)rus(ji) es,(h'_) (;) an(j) á(ji)li(hg)um(g_[uh:l]h) ex(i)spe(hi)ctá(g)mus?(g.) (:) Dí(j)ci(i_[uh:l]j)te(h) Jo(gf)án(gh)ni(hv_GF') quae(h) vi(ji)dí(g)stis :(g.) (:) ad(hj~) lu(j)men(i) réd(jk)e(hg)unt(g') cae(h)ci,(g.) (;) mór(h')tu(g)i(hg) re(f)súr(gh~)gunt,(h.) (;) páu(ji~)pe(j_k)res(kvJ'IH___) ev(ji)an(g)ge(hg)li(f)zán(gh~)tur,(h'_) (,) al(j_)le(h_jH'F)lú(hh){ia}.(g.) (::)}

  \newpage

  \begin{center}
    \textcolor{red}{\large{Oraison}}\\
  \end{center}

  \begin{multicols}{2}
    \parindent=0pt
    \begin{flushright}
      \textcolor{red}{\Vbar.} Dominus vobiscum.\\
      \textcolor{red}{\Rbar.} Et cum spiritu tuo.\\
    \end{flushright}

    \columnbreak
    
    \textit{\textcolor{red}{\Vbar.} Le Seigneur soit avec vous.\\
    \textcolor{red}{\Rbar.} Et avec votre esprit.}\\
  \end{multicols}

  \begin{multicols}{2}
    \parindent=0pt
    Excita, Dómine, corda nostra ad præparándas Unigéniti tui vias : \textcolor{red}{†} ut per ejus advéntum, \textcolor{red}{*} purificátis tibi méntibus servíre mereámur : \\ Qui tecum vivis et regnas in unitáte Spíritus Sancti Deus : per ómnia sæcula sæculórum. 
    \textcolor{red}{\Rbar.} Amen.

    \columnbreak

    \textit{ Réveillez, Seigneur, nos cœurs pour préparer une voie à votre Fils unique engendré, afin que, par sa venue, nous méritions de vous servir avec un esprit purifié. Lui qui avec vous vis et règne en l’unité du Saint Esprit, Dieu pour tous les siècles des siècles. 
    Amen.
    }
  \end{multicols}

  \medskip

  \begin{center}
    \textcolor{red}{\large{Conclusion de l'office}}
  \end{center}
  
  
  \begin{multicols}{2}
    \parindent=0pt
    \begin{flushright}
      \textcolor{red}{\Vbar.} Dominus vobiscum.\\
      \textcolor{red}{\Rbar.} Et cum spiritu tuo.\\
    \end{flushright}
  
    \columnbreak
    
    \textit{\textcolor{red}{\Vbar.} Le Seigneur soit avec vous.\\
    \textcolor{red}{\Rbar.} Et avec votre esprit.}\\
  \end{multicols}
  \bigskip
  \gresetinitiallines{1}
  \greillumination{\initfamily\fontsize{11mm}{11mm}\selectfont B}
  \gregorioscore{or--benedicamus_domino_(sundays_of_advent_and_lent)--solesmes_1961}
  \begin{center}
    \begin{footnotesize}
      \textcolor{red}{\textit{Sur un ton très grave : }}
    \end{footnotesize}
  \end{center}
  \begin{multicols}{2}
    \parindent=0pt
    \textcolor{red}{\Vbar.} Fidélium ánimæ per misericórdiam Dei requiéscant in pace.\\
    \textcolor{red}{\Rbar.} Amen.\\

    \columnbreak
    
    \textit{\textcolor{red}{\Vbar.} Que les âmes des fidèles défunts, par la
    miséricorde de Dieu, reposent en paix.\\
    \textcolor{red}{\Rbar.} Amen.}\\
  \end{multicols}

  \vspace*{\fill}\
  \begin{center}
    \greseparator{3}{30}
  \end{center}
  \vspace*{\fill}\

  \newpage

  \begin{center}
    \begin{LARGE}
      III\textsuperscript{e} dimanche de l'Avent
    \end{LARGE}
  \end{center}
  \medskip

  % ===== DEBUT Antienne =========
  \greillumination{\initfamily\fontsize{11mm}{11mm}\selectfont V}
  \gregorioscore{antiennes/an--veniet_dominus_et_non--solesmes_1961}
  \begin{center}
    \footnotesize{
      \textit{
        Le Seigneur viendra et il ne tardera pas, et il illuminera ce qui est caché dans les ténèbres, et il se manifestera à toutes les Nations, alléluia.
      }
    }
  \end{center}
  % \medskip
  % ===== FIN Antienne ===========

  % ===== DEBUT psaume ===========
  % gresetinitiallines : avec le parametre à 0, supprime l'ornement
  \gresetinitiallines{0}

  \begin{center}
    \normalsize{Psaume 109.}
  \end{center}

  \gregorioscore{psaumes/psaume109-Ia}
  \begin{enumerate}[label=\textcolor{red}{\emph{\arabic*}}]
    \setcounter{enumi}{2}
    \item Donec ponam ini\textbf{mí}cos \textbf{tu}os,~* scabéllum pe\textit{dum} \textit{tu}\textbf{ó}rum.

    \item Virgam virtútis tuæ emíttet Dómi\textbf{nus} ex \textbf{Si}on:~* domináre in médio inimicó\textit{rum} \textit{tu}\textbf{ó}rum.
    
    \item Tecum princípium in die virtútis tuæ in splendóri\textbf{bus} sanc\textbf{tó}rum:~* ex útero ante lucíferum \textit{gé}\textit{nu}\textbf{i} te.
    
    \item Jurávit Dóminus, et non pœni\textbf{té}bit \textbf{e}um:~* Tu es sacérdos in ætérnum secúndum órdi\textit{nem} \textit{Mel}\textbf{chí}sedech.
    
    \item Dóminus a \textbf{dex}tris \textbf{tu}is,~* confrégit in die iræ \textit{su}\textit{æ} \textbf{re}ges.
    
    \item Judicábit in natiónibus, im\textbf{plé}bit ru\textbf{í}nas:~* conquassábit cápita in ter\textit{ra} \textit{mul}\textbf{tó}rum.
    
    \item De torrénte in \textbf{vi}a \textbf{bi}bet:~* proptérea exal\textit{tá}\textit{bit} \textbf{ca}put.
    
    \item Glória \textbf{Pa}tri, et \textbf{Fí}lio,~* et Spirí\textit{tu}\textit{i} \textbf{Sanc}to.
    
    \item Sicut erat in princípio, et \textbf{nunc}, et \textbf{sem}per,~* et in sǽcula sæcu\textit{ló}\textit{rum}. \textbf{A}men.
  \end{enumerate}
  \medskip
  \grecommentary{\textit{Reprise de l'Antienne.}}
  \gabcsnippet{(c4) Ve(h')ni(g)et(g') Dó(h)mi(g)nus,(e'_[oh:h]) (,) et(f) non(g_[uh:l]h) tar(fe~)dá(d)bit,(d.) (;) ut(d) il(f)lú(g')mi(g)net(g) abs(g)cón(h')di(g)ta(f) te(g)ne(gh)brá(f)rum,(f.) (;) et(f) ma(g)ni(h')fe(g)stá(h')bit(g) se(e'_[oh:h]) (,) ad(f) o(g')mnes(f) gen(fe~)tes,(dc) al(d)le(fe)lú(d.){ia}.(d.) (::) }

  \newpage

  % ===== DEBUT Antienne =========
  \gresetinitiallines{1}
  \greillumination{\initfamily\fontsize{11mm}{11mm}\selectfont J}
  \gregorioscore{antiennes/an--jerusalem_gaude--solesmes_1961}
  \begin{center}
    \footnotesize{
      \textit{
        Jérusalem, réjouis-toi d’une grande joie, parce qu’un Sauveur viendra à toi, alléluia. 
      }
    }
  \end{center}
  % \medskip
  % ===== FIN Antienne ===========

  % ===== DEBUT psaume ===========
  % gresetinitiallines : avec le parametre à 0, supprime l'ornement
  \gresetinitiallines{0}

  \begin{center}
    \normalsize{Psaume 110.}
  \end{center}

  \gregorioscore{psaumes/psaume110-VIIb}

  \begin{enumerate}[label=\textcolor{red}{\emph{\arabic*}}]
    \setcounter{enumi}{1}
    \item Magna \textbf{ó}pera \textbf{Dó}mini:~* exquisíta in omnes volun\textbf{tá}tes \textbf{e}jus.

    \item Conféssio et magnificéntia \textbf{o}pus \textbf{e}jus:~* et justítia ejus manet in \textbf{sǽ}culum \textbf{sǽ}culi.
    
    \item Memóriam fecit mirabílium suórum,~† miséricors et mise\textbf{rá}tor \textbf{Dó}minus:~* escam dedit ti\textbf{mén}ti\textbf{bus} se.
    
    \item Memor erit in sǽculum testa\textbf{mén}ti \textbf{su}i:~* virtútem óperum suórum annuntiábit \textbf{pó}pulo \textbf{su}o:
    
    \item Ut det illis heredi\textbf{tá}tem \textbf{gén}tium:~* ópera mánuum ejus véritas, \textbf{et} ju\textbf{dí}cium.
    
    \item Fidélia ómnia mandáta ejus:~† confirmáta in \textbf{sǽ}culum \textbf{sǽ}culi,~* facta in veritáte et \textbf{æ}qui\textbf{tá}te.
    
    \item Redemptiónem misit \textbf{pó}pulo \textbf{su}o:~* mandávit in ætérnum testa\textbf{mén}tum \textbf{su}um.
    
    \item Sanctum, et terríbile \textbf{no}men \textbf{e}jus:~* inítium sapiéntiæ \textbf{ti}mor \textbf{Dó}mini.
    
    \item Intelléctus bonus ómnibus faci\textbf{én}tibus \textbf{e}um:~* laudátio ejus manet in \textbf{sǽ}culum \textbf{sǽ}culi.
    
    \item Glória \textbf{Pa}tri, et \textbf{Fí}lio,~* et Spi\textbf{rí}tui \textbf{Sanc}to.
    
    \item Sicut erat in princípio, et \textbf{nunc}, et \textbf{sem}per,~* et in sǽcula sæcu\textbf{ló}rum. \textbf{A}men.
  \end{enumerate}

  \medskip
  \grecommentary{\textit{Reprise de l'Antienne.}}
  \gabcsnippet{(c3) Je(f)rú(g')sa(h)lem(ii) gau(h_g)de(f'_) (,) gáu(h)di(g')o(h) ma(i_[uh:l]j)gno,(i.) (;) qui(j_l)a(k_[uh:l]l) vé(jk)ni(ji)et(i'_[oh:h]) ti(j)bi(i'_[oh:h]) (,) Sal(h)vá(hg)tor,(f) al(g)le(fg)lú(e.){ia}.(e.) (::)}

  \newpage

  % ===== DEBUT Antienne =========
  \gresetinitiallines{1}
  \greillumination{\initfamily\fontsize{11mm}{11mm}\selectfont D}
  \gregorioscore{antiennes/an--dabo_in_sion--solesmes_1961}
  \begin{center}
    \footnotesize{
      \textit{
        J’établirai dans Sion le salut, et dans Jérusalem ma gloire, alléluia.
      }
    }
  \end{center}
  % \medskip
  % ===== FIN Antienne ===========

  % ===== DEBUT psaume ===========
  % gresetinitiallines : avec le parametre à 0, supprime l'ornement

  \gresetinitiallines{0}

  \begin{center}
    \normalsize{Psaume 111.}\\
  \end{center}
  % \smallskip
  \grechangedim{baselineskip}{50pt}{scalable}
  \gregorioscore{psaumes/psaume111-VIIIG}
  \begin{enumerate}[label=\textcolor{red}{\emph{\arabic*}}]
    \setcounter{enumi}{1}
    \item Potens in terra erit semen \textbf{e}jus:~* generátio rectórum be\textit{ne}\textit{di}\textbf{cé}tur.

    \item Glória, et divítiæ in domo \textbf{e}jus:~* et justítia ejus manet in sǽ\textit{cu}\textit{lum} \textbf{sǽ}culi.
    
    \item Exórtum est in ténebris lumen \textbf{rec}tis:~* miséricors, et miserá\textit{tor}, \textit{et} \textbf{jus}tus.
    
    \item Jucúndus homo qui miserétur et cómmodat,~† dispónet sermónes suos in ju\textbf{dí}cio:~* quia in ætérnum non \textit{com}\textit{mo}\textbf{vé}bitur.
    
    \item In memória ætérna erit \textbf{jus}tus:~* ab auditióne mala \textit{non} \textit{ti}\textbf{mé}bit.
    
    \item Parátum cor ejus speráre in Dómino,~† confirmátum est cor \textbf{e}jus:~* non commovébitur donec despíciat ini\textit{mí}\textit{cos} \textbf{su}os.
    
    \item Dispérsit, dedit paupéribus:~† justítia ejus manet in sǽculum \textbf{sǽ}culi,~* cornu ejus exaltábi\textit{tur} \textit{in} \textbf{gló}ria.
    
    \item Peccátor vidébit, et irascétur,~† déntibus suis fremet et ta\textbf{bé}scet:~* desidérium peccató\textit{rum} \textit{per}\textbf{í}bit.
    
    \item Glória Patri, et \textbf{Fí}lio,~* et Spirí\textit{tu}\textit{i} \textbf{Sanc}to.
    
    \item Sicut erat in princípio, et nunc, et \textbf{sem}per,~* et in sǽcula sæcu\textit{ló}\textit{rum}. \textbf{A}men.
  \end{enumerate}

  \medskip
  \grecommentary{\textit{Reprise de l'Antienne.}}
  \gabcsnippet{(c4) Da(g)bo(gv_[oh:h]FD) in(fe~) Si(fg)on(ghhg) sa(fh)lú(h)tem,(g.) (;) et(h') in(j) Je(g)rú(hg)sa(fg)lem(g'_[oh:h]) (,) gló(h)ri(g')am(g) me(fe)am,(d) al(f)le(gh)lú(g.){ia}.(g.) (::)}

  \newpage

  % ===== DEBUT Antienne =========
  \gresetinitiallines{1}
  \greillumination{\initfamily\fontsize{11mm}{11mm}\selectfont M}
  \gregorioscore{antiennes/an--montes_et_omnes_colles--solesmes_1961}
  \begin{center}
    \footnotesize{\textit{
      Les montagnes et toutes les collines seront abaissées ; les chemins tortueux seront redressés, et les raboteux deviendront des voies aplanies : venez, Seigneur, et ne tardez pas, alléluia. 
    }}
  \end{center}

  % ===== FIN Antienne ===========

  % ===== DEBUT psaume ===========
  % gresetinitiallines : avec le parametre à 0, supprime l'ornement
  \gresetinitiallines{0}

  \begin{center}
    \normalsize{Psaume 112.}\\
  \end{center}
  % \smallskip

  \gregorioscore{psaumes/psaume112-Va}

  \begin{enumerate}[label=\textcolor{red}{\emph{\arabic*}}]
    \setcounter{enumi}{1}
    \item Sit nomen Dómini bene\textbf{díc}tum,~* ex hoc nunc, et \textbf{us}que in \textbf{sǽ}culum.

    \item A solis ortu usque ad oc\textbf{cá}sum,~* laudábile \textbf{no}men \textbf{Dó}mini.
    
    \item Excélsus super omnes gentes \textbf{Dó}minus,~* et super cælos \textbf{gló}ria \textbf{e}jus.
    
    \item Quis sicut Dóminus, Deus noster, qui in altis \textbf{há}bitat,~* et humília réspicit in cælo \textbf{et} in \textbf{ter}ra?
    
    \item Súscitans a terra \textbf{ín}opem,~* et de stércore \textbf{é}rigens \textbf{páu}perem:
    
    \item Ut cóllocet eum cum prin\textbf{cí}pibus,~* cum princípibus \textbf{pó}puli \textbf{su}i.
    
    \item Qui habitáre facit stérilem in \textbf{do}mo,~* matrem fili\textbf{ó}rum læ\textbf{tán}tem.
    
    \item Glória Patri, et \textbf{Fí}lio,~* et Spi\textbf{rí}tui \textbf{Sanc}to.
    
    \item Sicut erat in princípio, et nunc, et \textbf{sem}per,~* et in sǽcula sæcu\textbf{ló}rum. \textbf{A}men.
  \end{enumerate}

  \medskip
  \grecommentary{\textit{Reprise de l'Antienne.}}
  \gabcsnippet{(c3) Mon(f)tes(d) et(fh) o(hg)mnes(f) col(h_i)les(h.) (,) hu(h_i)mi(j)li(i')a(h)bún(hi)tur :(f.) (;) et(f) e(f')runt(i) pra(h)va(gh) in(f') di(e)ré(f)cta,(fef) (,) et(d_e) á(f')spe(e)ra(d') in(c) vi(e)as(ef) pla(d.)nas :(d.) (:) ve(h)ni(f) Dó(h_i)mi(h)ne,(gxggf.0) (;) et(d) no(e)li(e') tar(e)dá(e)re,(e') al(f)le(ef)lú(d.){ia}.(d.) (::)}

  \newpage

  % ===== DEBUT Antienne =========
  \gresetinitiallines{1}
  \greillumination{\initfamily\fontsize{11mm}{11mm}\selectfont J}
  \gregorioscore{antiennes/an--juste_et_pie_vivamus--solesmes}
  \begin{center}
    \footnotesize{
      \textit{
        Vivons justement et pieusement, attendant la bienheureuse espérance et l’avènement du Seigneur. 
      }
    }
  \end{center}
  % ===== DEBUT psaume ===========
  % gresetinitiallines : avec le parametre à 0, supprime l'ornement
  \begin{center}
    \normalsize{Psaume 113.}
  \end{center}

  % gresetinitiallines : avec le parametre à 0, supprime l'ornement
  \gresetinitiallines{0}
  \gregorioscore{psaumes/psaume113-IID}

  \begin{enumerate}[label=\textcolor{red}{\arabic*}]
    \setcounter{enumi}{1}
    \item Facta est Judǽa sanctificátio \textbf{e}jus,~* Israël potés\textit{tas} \textbf{e}jus.

    \item Mare vidit, et \textbf{fu}git:~* Jordánis convérsus est \textit{re}\textbf{trór}sum.
    
    \item Montes exsultavérunt ut a\textbf{rí}etes,~* et colles sicut a\textit{gni} \textbf{ó}vium.
    
    \item Quid est tibi, mare, quod fu\textbf{gís}ti:~* et tu, Jordánis, quia convérsus es \textit{re}\textbf{trór}sum?
    
    \item Montes, exsultástis sicut a\textbf{rí}etes,~* et colles, sicut a\textit{gni} \textbf{ó}vium.
    
    \item A fácie Dómini mota est \textbf{ter}ra,~* a fácie De\textit{i} \textbf{Ja}cob.
    
    \item Qui convértit petram in stagna a\textbf{quá}rum,~* et rupem in fontes \textit{a}\textbf{quá}rum.
    
    \item Non nobis, Dómine, non \textbf{no}bis:~* sed nómini tuo \textit{da} \textbf{gló}riam.
    
    \item Super misericórdia tua, et veritáte \textbf{tu}a:~* nequándo dicant gentes: Ubi est Deus \textit{e}\textbf{ó}rum?
    
    \item Deus autem noster in \textbf{cæ}lo:~* ómnia quæcúmque vólu\textit{it}, \textbf{fe}cit.
    
    \item Simulácra géntium argéntum, et \textbf{au}rum,~* ópera mánu\textit{um} \textbf{hó}minum.
    
    \item Os habent, et non lo\textbf{quén}tur:~* óculos habent, et non \textit{vi}\textbf{dé}bunt.
    
    \item Aures habent, et non \textbf{áu}dient:~* nares habent, et non o\textit{do}\textbf{rá}bunt.
    
    \item Manus habent, et non palpábunt:~† pedes habent, et non ambu\textbf{lá}bunt:~* non clamábunt in gúttu\textit{re} \textbf{su}o.
    
    \item Símiles illis fiant qui fáciunt \textbf{e}a:~* et omnes qui confídunt \textit{in} \textbf{e}is.
    
    \item Domus Israël sperávit in \textbf{Dó}mino:~* adjútor eórum et protéctor \textit{e}\textbf{ó}rum est,
    
    \item Domus Aaron sperávit in \textbf{Dó}mino:~* adjútor eórum et protéctor \textit{e}\textbf{ó}rum est,
    
    \item Qui timent Dóminum, speravérunt in \textbf{Dó}mino:~* adjútor eórum et protéctor \textit{e}\textbf{ó}rum est.
    
    \item Dóminus memor fuit \textbf{nos}tri:~* et benedí\textit{xit} \textbf{no}bis:
    
    \item Benedíxit dómui \textbf{Is}raël:~* benedíxit dómu\textit{i} \textbf{A}aron.
    
    \item Benedíxit ómnibus, qui timent \textbf{Dó}minum,~* pusíllis cum \textit{ma}\textbf{jó}ribus.
    
    \item Adjíciat Dóminus \textbf{su}per vos:~* super vos, et super fíli\textit{os} \textbf{ves}tros.
    
    \item Benedícti vos a \textbf{Dó}mino,~* qui fecit cælum, \textit{et} \textbf{ter}ram.
    
    \item Cælum cæli \textbf{Dó}mino:~* terram autem dedit fíli\textit{is} \textbf{hó}minum.
    
    \item Non mórtui laudábunt te, \textbf{Dó}mine:~* neque omnes, qui descéndunt in \textit{in}\textbf{fér}num.
    
    \item Sed nos qui vívimus, benedícimus \textbf{Dó}mino,~* ex hoc nunc et usque \textit{in} \textbf{sǽ}culum.
    
    \item Glória Patri, et \textbf{Fí}lio,~* et Spirítu\textit{i} \textbf{Sanc}to.
    
    \item Sicut erat in princípio, et nunc, et \textbf{sem}per,~* et in sǽcula sæculó\textit{rum}. \textbf{A}men.
  \end{enumerate}

  %  Répetition de l'Antienne
  \grecommentary{\textit{Reprise de l'Antienne.}}
  \gabcsnippet{(f3) Ju(ff)ste(e) et(f) pi(g_h)e(i) vi(hg)vá(f)mus,(f.) (;) ex(e)spe(f)ctán(h_/h_)tes(f') be(e)á(g)tam(ig~) spem,(h_g) (;) et(e) ad(f)vén(h')tum(g) Dó(e)mi(g)ni.(f.) (::)}

  \medskip
  \begin{center}
    \rule{2cm}{0.4pt}
  \end{center}
  \medskip

  \par \textit{\footnotesize\textcolor{red}{On se lève pour le Capitule.}}

  \begin{center}
    \textcolor{red}{\large{Capitule}}\\
    \small\textit{
      Épître aux Philippiens. IV, 4-5.
    }
  \end{center}

  \begin{multicols}{2}
    \parindent=0pt
    Fratres :  Gaudéte in Dómino semper : íterum dico, gau-dé-te. \textcolor{red}{†} Modéstia vestra nota sit ómni-bus ho-mí-nibus : \textcolor{red}{*} Dóminus enim prope est.  \\
    \textcolor{red}{\Rbar.} Deo grátias.

    \columnbreak

    \textit{ Frères : Réjouissez-vous dans le Seigneur en tout temps ; je le répète, réjouissez-vous. Que votre modération soit connue de tous les hommes : le Seigneur est proche. \\
    \textcolor{red}{\Rbar.} Rendons grâce à Dieu.
    }
  \end{multicols}


  \par \textit{\footnotesize\textcolor{red}{Le Célébrant entonne, ensuite, les Chantres et le Chœur alternent les versets. La Doxologie est chantée par tous.}}

  \begin{center}
    \textcolor{red}{\large{Hymne}}\\
  \end{center}
  
  \gresetinitiallines{1}
  \greillumination{\initfamily\fontsize{11mm}{11mm}\selectfont C}
  \gregorioscore{hymnes/hy--creator_alme_siderum--solesmes_1961}
  \bigskip

  \begin{center}
    \begin{footnotesize}
      \textcolor{red}{\textit{On chante le verset debout.}}
    \end{footnotesize}
    \begin{minipage}{0.8\linewidth}
      \gresetinitiallines{0}
      \gabcsnippet{(c3)<c><v>\Vbar</v>.</c> Ro(h)rá(h)te(h) cœ(h)li(h) dé(h)su(h)per(h), et(h) nu(h)bes(h) plu(h)ant(h) jus(h)tum. (g'_) (hvGF'Efgf.) (::) (Z) <c><v>\Rbar</v>.</c> Ape(h)riá(h)tur(h) te(h)rra(h) et(h) gér(h)mi(h)net(h) Sal(h)va(h)tó(h)rem. (g'_) (hvGF'Efgf.) (::)}
      \bigskip
      \begin{center}
        \textit{\textcolor{red}{\Vbar.} Cieux, répandez votre rosée, et vous, nuées, faites pleuvoir le juste.}\\
        \textit{\textcolor{red}{\Rbar.} Que la terre s’ouvre et qu’elle germe le Sauveur.}
      \end{center}
    \end{minipage}
  \end{center}
  \normalsize

  \vspace*{\fill}\
  \begin{center}
    \greseparator{3}{30}
  \end{center}
  \vspace*{\fill}\

  \newpage

  \begin{center}
    \textcolor{red}{\large{Antienne à Magnificat}}\\
    \begin{center}
      \footnotesize{\textit{
        A partir du 17 décembre, on prends l'antienne "O" correspondante. Voir page...
      }}
    \end{center}
  \end{center}

  \gresetinitiallines{1}
  \greillumination{\initfamily\fontsize{11mm}{11mm}\selectfont B}
  \gregorioscore{antiennes/an--beata_es_maria--solesmes_1934}
  \begin{center}
    \footnotesize{\textit{
      Bienheureuse êtes-vous, Marie, qui avez cru : car ce qui vous a été dit par le Seigneur s’accomplira en vous, alléluia. 
    }}
  \end{center}
  \medskip
  

  \gresetinitiallines{0}
  \gregorioscore{magnificat/magnificat-VIIIG}

  \begin{enumerate}[label=\textcolor{red}{\arabic*}]
    \setcounter{enumi}{2}
    \item Quia respéxit humilitátem ancíllæ \textbf{su}æ:~* ecce enim ex hoc beátam me dicent omnes gene\textit{ra}\textit{ti}\textbf{ó}nes.

    \item Quia fecit mihi magna qui \textbf{pot}ens est:~* et sanctum \textit{no}\textit{men} \textbf{e}jus.
    
    \item Et misericórdia ejus a progénie in pro\textbf{gé}nies~* timén\textit{ti}\textit{bus} \textbf{e}um.
    
    \item Fecit poténtiam in bráchio \textbf{su}o:~* dispérsit supérbos mente \textit{cor}\textit{dis} \textbf{su}i.
    
    \item Depósuit poténtes de \textbf{se}de,~* et exal\textit{tá}\textit{vit} \textbf{hú}miles.
    
    \item Esuriéntes implévit \textbf{bo}nis:~* et dívites dimí\textit{sit} \textit{in}\textbf{á}nes.
    
    \item Suscépit Israël púerum \textbf{su}um,~* recordátus misericór\textit{di}\textit{æ} \textbf{su}æ.
    
    \item Sicut locútus est ad patres \textbf{nos}tros,~* Abraham et sémini e\textit{jus} \textit{in} \textbf{sǽ}cula.
    \begin{center}
      \textit{\footnotesize \textcolor{red}{On attend la fin de l'encensement avant de chanter la doxologie.}}
    \end{center}
    \item Glória Patri, et \textbf{Fí}lio,~* et Spirí\textit{tu}\textit{i} \textbf{Sanc}to.
    
    \item Sicut erat in princípio, et nunc, et \textbf{sem}per,~* et in sǽcula sæcu\textit{ló}\textit{rum}. \textbf{A}men.
  \end{enumerate}

  % \grecommentary{\textit{Reprise de l'Antienne.}}
  \gabcsnippet{(c4) Be(g)á(gj)ta(hg) es(h') Ma(f)rí(g_h)a,(gh..) (,) quæ(gf) cre(g)di(h)dí(hk)sti :(jkj.) (;) per(h>')fi(i)ci(j')én(h)tur(h') in(g) te(g'_) (,) quæ(g) di(g_h)cta(f_e) sunt(d) ti(e_f)bi(gh) a(h') Dó(g)mi(fg)no,(gv_FD'_) al(f)le(fg)lú(g.)ia(g.) (::)}

  \newpage

  \begin{center}
    \textcolor{red}{\large{Oraison}}\\
  \end{center}

  \begin{multicols}{2}
    \parindent=0pt
    \begin{flushright}
      \textcolor{red}{\Vbar.} Dominus vobiscum.\\
      \textcolor{red}{\Rbar.} Et cum spiritu tuo.\\
    \end{flushright}

    \columnbreak
    
    \textit{\textcolor{red}{\Vbar.} Le Seigneur soit avec vous.\\
    \textcolor{red}{\Rbar.} Et avec votre esprit.}\\
  \end{multicols}

  \begin{multicols}{2}
    \parindent=0pt
    Aurem tuam, quæsumus, Dómine, \textcolor{red}{†} précibus nostris accómmoda : \textcolor{red}{*} et mentis nostræ ténebras grátia tuæ visitatiónis illústra : \\ Qui vivis et regnas cum Deo Patre, in unitáte Spíritus Sancti Deus : per ómnia sæcula sæculórum.
    \textcolor{red}{\Rbar.} Amen.

    \columnbreak

    \textit{ Prête ton oreille, Seigneur, à nos prières, et éclaire les ténèbres de notre esprit par la grâce de ta visite. Qui vivis et regnas cum Deo Patre, in unitáte Spíritus Sancti Deus : per ómnia sæcula sæculórum. Toi qui vis et règne avec Dieu le Père en l’unité du Saint Esprit, Dieu pour tous les siècles des siècles.
    Amen.
    }
  \end{multicols}

  \medskip

  \begin{center}
    \textcolor{red}{\large{Conclusion de l'office}}
  \end{center}
  
  
  \begin{multicols}{2}
    \parindent=0pt
    \begin{flushright}
      \textcolor{red}{\Vbar.} Dominus vobiscum.\\
      \textcolor{red}{\Rbar.} Et cum spiritu tuo.\\
    \end{flushright}
  
    \columnbreak
    
    \textit{\textcolor{red}{\Vbar.} Le Seigneur soit avec vous.\\
    \textcolor{red}{\Rbar.} Et avec votre esprit.}\\
  \end{multicols}
  \bigskip
  \gresetinitiallines{1}
  \greillumination{\initfamily\fontsize{11mm}{11mm}\selectfont B}
  \gregorioscore{or--benedicamus_domino_(sundays_of_advent_and_lent)--solesmes_1961}
  \begin{center}
    \begin{footnotesize}
      \textcolor{red}{\textit{Sur un ton très grave : }}
    \end{footnotesize}
  \end{center}
  \begin{multicols}{2}
    \parindent=0pt
    \textcolor{red}{\Vbar.} Fidélium ánimæ per misericórdiam Dei requiéscant in pace.\\
    \textcolor{red}{\Rbar.} Amen.\\

    \columnbreak
    
    \textit{\textcolor{red}{\Vbar.} Que les âmes des fidèles défunts, par la
    miséricorde de Dieu, reposent en paix.\\
    \textcolor{red}{\Rbar.} Amen.}\\
  \end{multicols}

  \vspace*{\fill}\
  \begin{center}
    \greseparator{3}{30}
  \end{center}
  \vspace*{\fill}\

  \newpage

  \begin{center}
    \begin{LARGE}
      IV\textsuperscript{e} dimanche de l'Avent
    \end{LARGE}
  \end{center}
  \medskip

  % ===== DEBUT Antienne =========
  \greillumination{\initfamily\fontsize{11mm}{11mm}\selectfont C}
  \gregorioscore{antiennes/an--canite_tuba--solesmes_1961}
  \begin{center}
    \footnotesize{
      \textit{
        Sonnez de la trompette dans Sion, parce que le jour du Seigneur est proche : voici qu’il vient pour nous sauver, alléluia, alléluia. 
      }
    }
  \end{center}
  % \medskip
  % ===== FIN Antienne ===========

  % ===== DEBUT psaume ===========
  % gresetinitiallines : avec le parametre à 0, supprime l'ornement
  \gresetinitiallines{0}

  \begin{center}
    \normalsize{Psaume 109.}
  \end{center}

  \gregorioscore{psaumes/psaume109-Ig}
  \begin{enumerate}[label=\textcolor{red}{\emph{\arabic*}}]
    \setcounter{enumi}{2}
    \item Donec ponam ini\textbf{mí}cos \textbf{tu}os,~* scabéllum pe\textit{dum} \textit{tu}\textbf{ó}rum.

    \item Virgam virtútis tuæ emíttet Dómi\textbf{nus} ex \textbf{Si}on:~* domináre in médio inimicó\textit{rum} \textit{tu}\textbf{ó}rum.

    \item Tecum princípium in die virtútis tuæ in splendóri\textbf{bus} sanc\textbf{tó}rum:~* ex útero ante lucíferum \textit{gé}\textit{nu}\textbf{i} te.

    \item Jurávit Dóminus, et non pœni\textbf{té}bit \textbf{e}um:~* Tu es sacérdos in ætérnum secúndum órdi\textit{nem} \textit{Mel}\textbf{chí}sedech.

    \item Dóminus a \textbf{dex}tris \textbf{tu}is,~* confrégit in die iræ \textit{su}\textit{æ} \textbf{re}ges.

    \item Judicábit in natiónibus, im\textbf{plé}bit ru\textbf{í}nas:~* conquassábit cápita in ter\textit{ra} \textit{mul}\textbf{tó}rum.

    \item De torrénte in \textbf{vi}a \textbf{bi}bet:~* proptérea exal\textit{tá}\textit{bit} \textbf{ca}put.

    \item Glória \textbf{Pa}tri, et \textbf{Fí}lio,~* et Spirí\textit{tu}\textit{i} \textbf{Sanc}to.

    \item Sicut erat in princípio, et \textbf{nunc}, et \textbf{sem}per,~* et in sǽcula sæcu\textit{ló}\textit{rum}. \textbf{A}men.
  \end{enumerate}
  \medskip
  \grecommentary{\textit{Reprise de l'Antienne.}}
  \gabcsnippet{(c4) Ca(f)ni(c')te(d) tu(ixdh'!iv)ba(h'_) (,) in(h_ggf~) Si(ixghg/h_i)on,(h_gg) (;) qui(f)a(g) pro(h')pe(g) est(f') di(g)es(d_g) Dó(gvF'ED'___ew!fv_ED)mi(c!ew!fv_ED'e)ni :(ed..) (:) ec(f)ce(f) vé(f_g)ni(f)et(c') ad(d) sal(f)ván(e)dum(c) nos,(d!ewf.) (;) al(f)le(ed)lú(ffg){ia},(f_e) al(g)le(fe)lú(d.){ia}.(d.) (::)}

  \newpage

  % ===== DEBUT Antienne =========
  \gresetinitiallines{1}
  \greillumination{\initfamily\fontsize{11mm}{11mm}\selectfont E}
  \gregorioscore{antiennes/an--ecce_veniet_desideratus--solesmes_1961}
  \begin{center}
    \footnotesize{
      \textit{
        Voici que viendra le désiré de toutes les Nations, et la maison du Seigneur sera remplie de gloire, alléluia.  
      }
    }
  \end{center}
  % \medskip
  % ===== FIN Antienne ===========

  % ===== DEBUT psaume ===========
  % gresetinitiallines : avec le parametre à 0, supprime l'ornement
  \gresetinitiallines{0}

  \begin{center}
    \normalsize{Psaume 110.}
  \end{center}

  \gregorioscore{psaumes/psaume110-If}

  \begin{enumerate}[label=\textcolor{red}{\emph{\arabic*}}]
    \setcounter{enumi}{1}
    \item Magna \textbf{ó}pera \textbf{Dó}mini:~* exquisíta in omnes volun\textit{tá}\textit{tes} \textbf{e}jus.

    \item Conféssio et magnificéntia \textbf{o}pus \textbf{e}jus:~* et justítia ejus manet in sǽ\textit{cu}\textit{lum} \textbf{sǽ}culi.
    
    \item Memóriam fecit mirabílium suórum,~† miséricors et mise\textbf{rá}tor \textbf{Dó}minus:~* escam dedit ti\textit{mén}\textit{ti}\textbf{bus} se.
    
    \item Memor erit in sǽculum testa\textbf{mén}ti \textbf{su}i:~* virtútem óperum suórum annuntiábit pó\textit{pu}\textit{lo} \textbf{su}o:
    
    \item Ut det illis heredi\textbf{tá}tem \textbf{gén}tium:~* ópera mánuum ejus véritas, \textit{et} \textit{ju}\textbf{dí}cium.
    
    \item Fidélia ómnia mandáta ejus:~† confirmáta in \textbf{sǽ}culum \textbf{sǽ}culi,~* facta in veritáte et \textit{æ}\textit{qui}\textbf{tá}te.
    
    \item Redemptiónem misit \textbf{pó}pulo \textbf{su}o:~* mandávit in ætérnum testa\textit{mén}\textit{tum} \textbf{su}um.
    
    \item Sanctum, et terríbile \textbf{no}men \textbf{e}jus:~* inítium sapiéntiæ \textit{ti}\textit{mor} \textbf{Dó}mini.
    
    \item Intelléctus bonus ómnibus faci\textbf{én}tibus \textbf{e}um:~* laudátio ejus manet in sǽ\textit{cu}\textit{lum} \textbf{sǽ}culi.
    
    \item Glória \textbf{Pa}tri, et \textbf{Fí}lio,~* et Spirí\textit{tu}\textit{i} \textbf{Sanc}to.
    
    \item Sicut erat in princípio, et \textbf{nunc}, et \textbf{sem}per,~* et in sǽcula sæcu\textit{ló}\textit{rum}. \textbf{A}men.
  \end{enumerate}

  \medskip
  \grecommentary{\textit{Reprise de l'Antienne.}}
  \gabcsnippet{(c4) Ec(d)ce(d) vé(d)ni(d)et(d) <c>*</c>() de(dc)si(f)de(g)rá(f_h)tus(h.) (,) cun(ixi_[oh:h]h)ctis(g.) gén(f)ti(gh)bus :(g.) (;) et(g) re(g)plé(g')bi(e)tur(g') gló(h)ri(gf)a(f'_) (,) do(g)mus(e') Dó(f)mi(e)ni,(dc) al(d)le(fe)lú(d.){ia}.(d.) (::)}

  \newpage

  % ===== DEBUT Antienne =========
  \gresetinitiallines{1}
  \greillumination{\initfamily\fontsize{11mm}{11mm}\selectfont E}
  \gregorioscore{antiennes/an--erunt_prava--solesmes_1961}
  \begin{center}
    \footnotesize{
      \textit{
        Les chemins tortueux seront redressés et les raboteux deviendront des voies aplanies : venez, Seigneur, et ne tardez pas, alléluia.
      }
    }
  \end{center}
  % \medskip
  % ===== FIN Antienne ===========

  % ===== DEBUT psaume ===========
  % gresetinitiallines : avec le parametre à 0, supprime l'ornement

  \gresetinitiallines{0}

  \begin{center}
    \normalsize{Psaume 111.}\\
  \end{center}
  % \smallskip
  \grechangedim{baselineskip}{50pt}{scalable}
  \gregorioscore{psaumes/psaume111-Ig}
  \begin{enumerate}[label=\textcolor{red}{\emph{\arabic*}}]
    \setcounter{enumi}{1}
    \item Potens in terra erit \textbf{se}men \textbf{e}jus:~* generátio rectórum be\textit{ne}\textit{di}\textbf{cé}tur.

    \item Glória, et divítiæ in \textbf{do}mo \textbf{e}jus:~* et justítia ejus manet in sǽ\textit{cu}\textit{lum} \textbf{sǽ}culi.

    \item Exórtum est in ténebris \textbf{lu}men \textbf{rec}tis:~* miséricors, et miserá\textit{tor}, \textit{et} \textbf{jus}tus.

    \item Jucúndus homo qui miserétur et cómmodat,~† dispónet sermónes suos \textbf{in} ju\textbf{dí}cio:~* quia in ætérnum non \textit{com}\textit{mo}\textbf{vé}bitur.

    \item In memória ætérna \textbf{e}rit \textbf{jus}tus:~* ab auditióne mala \textit{non} \textit{ti}\textbf{mé}bit.

    \item Parátum cor ejus speráre in Dómino,~† confirmátum \textbf{est} cor \textbf{e}jus:~* non commovébitur donec despíciat ini\textit{mí}\textit{cos} \textbf{su}os.

    \item Dispérsit, dedit paupéribus:~† justítia ejus manet in \textbf{sǽ}culum \textbf{sǽ}culi,~* cornu ejus exaltábi\textit{tur} \textit{in} \textbf{gló}ria.

    \item Peccátor vidébit, et irascétur,~† déntibus suis fremet \textbf{et} ta\textbf{bé}scet:~* desidérium peccató\textit{rum} \textit{per}\textbf{í}bit.

    \item Glória \textbf{Pa}tri, et \textbf{Fí}lio,~* et Spirí\textit{tu}\textit{i} \textbf{Sanc}to.

    \item Sicut erat in princípio, et \textbf{nunc}, et \textbf{sem}per,~* et in sǽcula sæcu\textit{ló}\textit{rum}. \textbf{A}men.
  \end{enumerate}

  \medskip
  \grecommentary{\textit{Reprise de l'Antienne.}}
  \gabcsnippet{(c4) E(c)runt(d) pra(ixdh)va(h_i)() in(h) di(g')ré(h)cta,(hgh) (,) et(f_g) á(h')spe(g)ra(f') in(e) vi(g)as(fe) pla(d.)nas :(d.) (:) ve(f)ni(g) Dó(e)mi(f)ne,(c.) (,) et(e) no(g)li(f') tar(f)dá(fe)re,(dc) al(d)le(fe)lú(d.){ia}.(d.) (::)}

  \newpage

  % ===== DEBUT Antienne =========
  \gresetinitiallines{1}
  \greillumination{\initfamily\fontsize{11mm}{11mm}\selectfont D}
  \gregorioscore{antiennes/an--dominus_veniet--solesmes_1961}
  \begin{center}
    \footnotesize{\textit{
      Le Seigneur va venir, allez au-devant de lui, disant : Sa puissance est grande et son règne n’aura pas de fin ; il est Dieu, Fort, Dominateur, Prince de la paix, alléluia, alléluia. 
    }}
  \end{center}

  % ===== FIN Antienne ===========

  % ===== DEBUT psaume ===========
  % gresetinitiallines : avec le parametre à 0, supprime l'ornement
  \gresetinitiallines{0}

  \begin{center}
    \normalsize{Psaume 112.}\\
  \end{center}
  % \smallskip

  \gregorioscore{psaumes/psaume112-If}

  \begin{enumerate}[label=\textcolor{red}{\emph{\arabic*}}]
    \setcounter{enumi}{1}
    \item Sit nomen Dómini \textbf{be}ne\textbf{díc}tum,~* ex hoc nunc, et us\textit{que} \textit{in} \textbf{sǽ}culum.

    \item A solis ortu usque \textbf{ad} oc\textbf{cá}sum,~* laudábile \textit{no}\textit{men} \textbf{Dó}mini.
    
    \item Excélsus super omnes \textbf{gen}tes \textbf{Dó}minus,~* et super cælos gló\textit{ri}\textit{a} \textbf{e}jus.
    
    \item Quis sicut Dóminus, Deus noster, qui in \textbf{al}tis \textbf{há}bitat,~* et humília réspicit in cælo \textit{et} \textit{in} \textbf{ter}ra?
    
    \item Súscitans a \textbf{ter}ra \textbf{ín}opem,~* et de stércore é\textit{ri}\textit{gens} \textbf{páu}perem:
    
    \item Ut cóllocet eum \textbf{cum} prin\textbf{cí}pibus,~* cum princípibus pó\textit{pu}\textit{li} \textbf{su}i.
    
    \item Qui habitáre facit stéri\textbf{lem} in \textbf{do}mo,~* matrem filió\textit{rum} \textit{læ}\textbf{tán}tem.
    
    \item Glória \textbf{Pa}tri, et \textbf{Fí}lio,~* et Spirí\textit{tu}\textit{i} \textbf{Sanc}to.
    
    \item Sicut erat in princípio, et \textbf{nunc}, et \textbf{sem}per,~* et in sǽcula sæcu\textit{ló}\textit{rum}. \textbf{A}men.
  \end{enumerate}

  \medskip
  \grecommentary{\textit{Reprise de l'Antienne.}}
  \gabcsnippet{(c4) Do(d_c)mi(f)nus(g') vé(f)ni(gh)et,(h'_) (,) oc(h)cúr(hj)ri(h)te(g') il(i)li,(g') di(e)cén(ff)tes :(d.) (;) Ma(h_ih)gnum(g) prin(hj~)cí(j)pi(jv_IH)um,(h.) (;) et(h_gg) re(f)gni(gh) e(g_[oh:h]fg)jus(g') non(h) e(g)rit(f') fi(g)nis :(d.) (:) De(g_2[oh:h]g_2[oh:h]g_[oh:h])us,(g_[oh:h]d) for(f)tis,(d) do(e)mi(f)ná(d_e)tor,(ddc.) (;) prin(c_[uh:l-0.5mm]d)ceps(f_e) pa(f_g)cis,(gg) (,) al(c)le(d')lú(f){ia},(e') al(g)le(fe)lú(d.){ia}.(d.) (::)}

  \newpage

  % ===== DEBUT Antienne =========
  \gresetinitiallines{1}
  \greillumination{\initfamily\fontsize{11mm}{11mm}\selectfont O}
  \gregorioscore{antiennes/an--omnipotens_sermo--solesmes_1961}
  \begin{center}
    \footnotesize{
      \textit{
        Votre parole toute-puissante, Seigneur, viendra du trône royal, alléluia. 
      }
    }
  \end{center}
  % ===== DEBUT psaume ===========
  % gresetinitiallines : avec le parametre à 0, supprime l'ornement
  \begin{center}
    \normalsize{Psaume 113.}
  \end{center}

  % gresetinitiallines : avec le parametre à 0, supprime l'ornement
  \gresetinitiallines{0}
  \gregorioscore{psaumes/psaume113-IID}

  \begin{enumerate}[label=\textcolor{red}{\arabic*}]
    \setcounter{enumi}{1}
    \item Facta est Judǽa sanctificátio \textbf{e}jus,~* Israël potés\textit{tas} \textbf{e}jus.

    \item Mare vidit, et \textbf{fu}git:~* Jordánis convérsus est \textit{re}\textbf{trór}sum.
    
    \item Montes exsultavérunt ut a\textbf{rí}etes,~* et colles sicut a\textit{gni} \textbf{ó}vium.
    
    \item Quid est tibi, mare, quod fu\textbf{gís}ti:~* et tu, Jordánis, quia convérsus es \textit{re}\textbf{trór}sum?
    
    \item Montes, exsultástis sicut a\textbf{rí}etes,~* et colles, sicut a\textit{gni} \textbf{ó}vium.
    
    \item A fácie Dómini mota est \textbf{ter}ra,~* a fácie De\textit{i} \textbf{Ja}cob.
    
    \item Qui convértit petram in stagna a\textbf{quá}rum,~* et rupem in fontes \textit{a}\textbf{quá}rum.
    
    \item Non nobis, Dómine, non \textbf{no}bis:~* sed nómini tuo \textit{da} \textbf{gló}riam.
    
    \item Super misericórdia tua, et veritáte \textbf{tu}a:~* nequándo dicant gentes: Ubi est Deus \textit{e}\textbf{ó}rum?
    
    \item Deus autem noster in \textbf{cæ}lo:~* ómnia quæcúmque vólu\textit{it}, \textbf{fe}cit.
    
    \item Simulácra géntium argéntum, et \textbf{au}rum,~* ópera mánu\textit{um} \textbf{hó}minum.
    
    \item Os habent, et non lo\textbf{quén}tur:~* óculos habent, et non \textit{vi}\textbf{dé}bunt.
    
    \item Aures habent, et non \textbf{áu}dient:~* nares habent, et non o\textit{do}\textbf{rá}bunt.
    
    \item Manus habent, et non palpábunt:~† pedes habent, et non ambu\textbf{lá}bunt:~* non clamábunt in gúttu\textit{re} \textbf{su}o.
    
    \item Símiles illis fiant qui fáciunt \textbf{e}a:~* et omnes qui confídunt \textit{in} \textbf{e}is.
    
    \item Domus Israël sperávit in \textbf{Dó}mino:~* adjútor eórum et protéctor \textit{e}\textbf{ó}rum est,
    
    \item Domus Aaron sperávit in \textbf{Dó}mino:~* adjútor eórum et protéctor \textit{e}\textbf{ó}rum est,
    
    \item Qui timent Dóminum, speravérunt in \textbf{Dó}mino:~* adjútor eórum et protéctor \textit{e}\textbf{ó}rum est.
    
    \item Dóminus memor fuit \textbf{nos}tri:~* et benedí\textit{xit} \textbf{no}bis:
    
    \item Benedíxit dómui \textbf{Is}raël:~* benedíxit dómu\textit{i} \textbf{A}aron.
    
    \item Benedíxit ómnibus, qui timent \textbf{Dó}minum,~* pusíllis cum \textit{ma}\textbf{jó}ribus.
    
    \item Adjíciat Dóminus \textbf{su}per vos:~* super vos, et super fíli\textit{os} \textbf{ves}tros.
    
    \item Benedícti vos a \textbf{Dó}mino,~* qui fecit cælum, \textit{et} \textbf{ter}ram.
    
    \item Cælum cæli \textbf{Dó}mino:~* terram autem dedit fíli\textit{is} \textbf{hó}minum.
    
    \item Non mórtui laudábunt te, \textbf{Dó}mine:~* neque omnes, qui descéndunt in \textit{in}\textbf{fér}num.
    
    \item Sed nos qui vívimus, benedícimus \textbf{Dó}mino,~* ex hoc nunc et usque \textit{in} \textbf{sǽ}culum.
    
    \item Glória Patri, et \textbf{Fí}lio,~* et Spirítu\textit{i} \textbf{Sanc}to.
    
    \item Sicut erat in princípio, et nunc, et \textbf{sem}per,~* et in sǽcula sæculó\textit{rum}. \textbf{A}men.
  \end{enumerate}

  %  Répetition de l'Antienne
  \grecommentary{\textit{Reprise de l'Antienne.}}
  \gabcsnippet{(f3) O(h)mní(h)pot(gh)ens(f') (,) ser(f)mo(e) tu(f')us(h) Dó(h')mi(g)ne(e.) (;) a(g) re(h)gá(i'_[oh:h])li(h)bus(f) sé(hh)di(g')bus(h) vé(i)ni(h)et,(g_[oh:h]e) al(f)le(hg)lú(f.){ia}.(f.) (::)}

  \medskip
  \begin{center}
    \rule{2cm}{0.4pt}
  \end{center}
  \medskip

  \par \textit{\footnotesize\textcolor{red}{On se lève pour le Capitule.}}

  \begin{center}
    \textcolor{red}{\large{Capitule}}\\
    \small\textit{
      I\textsuperscript{ère.} Épître aux Corinthiens. IV.
    }
  \end{center}

  \begin{multicols}{2}
    \parindent=0pt
    Fratres :  Sic nos exístimet homo ut minístros Christi, et dispensatóres mysteriórum Dei. \textcolor{red}{†} Hic jam quæritur inter dispensatóres,  \textcolor{red}{*} ut fidélis quis inveniátur.   \\
    \textcolor{red}{\Rbar.} Deo grátias.

    \columnbreak

    \textit{ Frères : Ainsi, qu’on nous regarde comme des serviteurs du Christ et des dispensateurs des mystères de Dieu. Eh bien ! ce que l’on cherche dans les dispensateurs, c’est que chacun soit trouvé fidèle. \\
    \textcolor{red}{\Rbar.} Rendons grâce à Dieu.
    }
  \end{multicols}


  \par \textit{\footnotesize\textcolor{red}{Le Célébrant entonne, ensuite, les Chantres et le Chœur alternent les versets. La Doxologie est chantée par tous.}}

  \begin{center}
    \textcolor{red}{\large{Hymne}}\\
  \end{center}
  
  \gresetinitiallines{1}
  \greillumination{\initfamily\fontsize{11mm}{11mm}\selectfont C}
  \gregorioscore{hymnes/hy--creator_alme_siderum--solesmes_1961}
  \bigskip

  \begin{center}
    \begin{footnotesize}
      \textcolor{red}{\textit{On chante le verset debout.}}
    \end{footnotesize}
    \begin{minipage}{0.8\linewidth}
      \gresetinitiallines{0}
      \gabcsnippet{(c3)<c><v>\Vbar</v>.</c> Ro(h)rá(h)te(h) cœ(h)li(h) dé(h)su(h)per(h), et(h) nu(h)bes(h) plu(h)ant(h) jus(h)tum. (g'_) (hvGF'Efgf.) (::) (Z) <c><v>\Rbar</v>.</c> Ape(h)riá(h)tur(h) te(h)rra(h) et(h) gér(h)mi(h)net(h) Sal(h)va(h)tó(h)rem. (g'_) (hvGF'Efgf.) (::)}
      \bigskip
      \begin{center}
        \textit{\textcolor{red}{\Vbar.} Cieux, répandez votre rosée, et vous, nuées, faites pleuvoir le juste.}\\
        \textit{\textcolor{red}{\Rbar.} Que la terre s’ouvre et qu’elle germe le Sauveur.}
      \end{center}
    \end{minipage}
  \end{center}
  \normalsize

  \vspace*{\fill}\
  \begin{center}
    \greseparator{3}{30}
  \end{center}
  \vspace*{\fill}\

  \newpage

  \begin{center}
    \textcolor{red}{\large{Antienne à Magnificat}}\\
    \begin{center}
      \footnotesize{\textit{
        Prendre une des antiennes "O" page ...
      }}
    \end{center}
  \end{center}


  \begin{center}
    \textcolor{red}{\large{Oraison}}\\
  \end{center}

  \begin{multicols}{2}
    \parindent=0pt
    \begin{flushright}
      \textcolor{red}{\Vbar.} Dominus vobiscum.\\
      \textcolor{red}{\Rbar.} Et cum spiritu tuo.\\
    \end{flushright}

    \columnbreak
    
    \textit{\textcolor{red}{\Vbar.} Le Seigneur soit avec vous.\\
    \textcolor{red}{\Rbar.} Et avec votre esprit.}\\
  \end{multicols}

  \begin{multicols}{2}
    \parindent=0pt
    Excita, quæsumus, Dómine, poténtiam tuam, et veni, et magna nobis virtúte succúrre ; \textcolor{red}{†} ut per auxílium grátiæ tuæ, quod nostra peccáta præpédiunt, \textcolor{red}{*} indulgéntia tuæ propitiatiónis accéleret : \\
    Qui vivis et regnas cum Deo Patre, in unitáte Spíritus Sancti Deus, per ómnia sæcula sæculórum. 
    \textcolor{red}{\Rbar.} Amen.

    \columnbreak

    \textit{ Faites paraître, Seigneur, votre puissance, et venez ; dans votre grande force, secourez-nous, afin que, par le secours de votre grâce, votre indulgence se hâte de nous donner le bienfait que nos péchés nous empêchent d’obtenir. Vous qui vivez et règnez avec Dieu le Père en l’unité du Saint Esprit, Dieu pour tous les siècles des siècles. 
    Amen.
    }
  \end{multicols}

  \medskip

  \begin{center}
    \textcolor{red}{\large{Conclusion de l'office}}
  \end{center}
  
  
  \begin{multicols}{2}
    \parindent=0pt
    \begin{flushright}
      \textcolor{red}{\Vbar.} Dominus vobiscum.\\
      \textcolor{red}{\Rbar.} Et cum spiritu tuo.\\
    \end{flushright}
  
    \columnbreak
    
    \textit{\textcolor{red}{\Vbar.} Le Seigneur soit avec vous.\\
    \textcolor{red}{\Rbar.} Et avec votre esprit.}\\
  \end{multicols}
  \bigskip
  \gresetinitiallines{1}
  \greillumination{\initfamily\fontsize{11mm}{11mm}\selectfont B}
  \gregorioscore{or--benedicamus_domino_(sundays_of_advent_and_lent)--solesmes_1961}
  \begin{center}
    \begin{footnotesize}
      \textcolor{red}{\textit{Sur un ton très grave : }}
    \end{footnotesize}
  \end{center}
  \begin{multicols}{2}
    \parindent=0pt
    \textcolor{red}{\Vbar.} Fidélium ánimæ per misericórdiam Dei requiéscant in pace.\\
    \textcolor{red}{\Rbar.} Amen.\\

    \columnbreak
    
    \textit{\textcolor{red}{\Vbar.} Que les âmes des fidèles défunts, par la
    miséricorde de Dieu, reposent en paix.\\
    \textcolor{red}{\Rbar.} Amen.}\\
  \end{multicols}

  \begin{center}
    \textcolor{red}{\large{Salut du Saint Sacrement page ...}}
  \end{center}

  \vspace*{\fill}\
  \begin{center}
    \greseparator{3}{30}
  \end{center}
  \vspace*{\fill}\

  \newpage

  \begin{center}
    \begin{LARGE}
      Les grandes antiennes "O"
    \end{LARGE}
  \end{center}
  \medskip
  
  \begin{center}
    \textcolor{red}{\large{Le 17 décembre}}
  \end{center}

  \gresetinitiallines{1}
  \greillumination{\initfamily\fontsize{11mm}{11mm}\selectfont O}
  \gregorioscore{antiennes/an--o_sapientia--solesmes_1961}
  \begin{center}
    \footnotesize{
      \textit{
        O Sagesse, qui êtes sortie de la bouche du Très-Haut (Ecclésiastique, XXIV, 3), atteignant d’une extrémité à une autre extrémité, et disposant toutes choses avec force et douceur (Sagesse VIII, 1) : venez pour nous enseigner la voie de la prudence. 
      }
    }
  \end{center}
  \medskip

  \gresetinitiallines{0}
  \gregorioscore{magnificat/magnificat-IIDsolemn}

  \begin{enumerate}[label=\textcolor{red}{\arabic*}]
    \setcounter{enumi}{2}
    \item Quia respéxit humilitátem \textit{an}\textit{cíl}\textit{læ} \textbf{su}æ:~* ecce enim ex hoc beátam me dicent omnes genera\textit{ti}\textbf{ó}nes.

    \item Quia fecit mihi \textit{ma}\textit{gna} \textit{qui} \textbf{pot}ens est:~* et sanctum no\textit{men} \textbf{e}jus.
    
    \item Et misericórdia ejus a progéni\textit{e} \textit{in} \textit{pro}\textbf{gé}nies~* timénti\textit{bus} \textbf{e}um.
    
    \item Fecit poténtiam in \textit{brá}\textit{chi}\textit{o} \textbf{su}o:~* dispérsit supérbos mente cor\textit{dis} \textbf{su}i.
    
    \item Depósuit pot\textit{én}\textit{tes} \textit{de} \textbf{se}de,~* et exaltá\textit{vit} \textbf{hú}miles.
    
    \item Esuriéntes \textit{im}\textit{plé}\textit{vit} \textbf{bo}nis:~* et dívites dimísit \textit{in}\textbf{á}nes.
    
    \item Suscépit Israël \textit{pú}\textit{e}\textit{rum} \textbf{su}um,~* recordátus misericórdi\textit{æ} \textbf{su}æ.
    
    \item Sicut locútus est \textit{ad} \textit{pa}\textit{tres} \textbf{nos}tros,~* Abraham et sémini ejus \textit{in} \textbf{sǽ}cula.
    
    \item Glória \textit{Pa}\textit{tri}, \textit{et} \textbf{Fí}lio,~* et Spirítu\textit{i} \textbf{Sanc}to.
    
    \item Sicut erat in princípio, \textit{et} \textit{nunc}, \textit{et} \textbf{sem}per,~* et in sǽcula sæculó\textit{rum}. \textbf{A}men.
  \end{enumerate}

  \grecommentary{\textit{Reprise de l'Antienne.}}
  \gabcsnippet{(f3) O(ehhg) Sa(hg)pi(f)én(gf~)ti(ef)a,(f'_) (,) quae(f) ex(f) o(fg)re(f) Al(f)tís(g)si(f)mi(g_[uh:l]h) pro(fg)dí(f)sti,(e.) (;) at(f)tín(cf~)gens(f) a(f) fi(g)ne(f') us(g)que(f) ad(f) fi(g')nem,(f) fór(g')ti(h)ter(kxiji/jkj.) (;) su(j)á(j')vi(i)ter(h) dis(g)po(h)néns(ih~)que(f') ó(g)mni(fe)a :(e.) (:) ve(eh)ni(g'_[oh:h]) (,) ad(h) do(f)cén(gf~)dum(ef~) nos(fv_EC.) (,) vi(e_[uh:l]f)am(h) pru(g)dén(e)ti(g)ae.(f.) (::)}

  \newpage

  \begin{center}
    \textcolor{red}{\large{Le 18 décembre}}
  \end{center}

  \gresetinitiallines{1}
  \greillumination{\initfamily\fontsize{11mm}{11mm}\selectfont O}
  \gregorioscore{antiennes/an--o_adonai--solesmes_1961}
  \begin{center}
    \footnotesize{
      \textit{
        O Adonaï, et Conducteur de la maison d’Israël (Exode VI, 2-3, 13), qui avez apparu à Moïse dans le feu du buisson ardent (Exode III, 2) et lui avez donné la loi sur le Sinaï : venez pour nous racheter par la puissance de votre bras. (Exode VI, 6) 
      }
    }
  \end{center}
  \medskip

  \gresetinitiallines{0}
  \gregorioscore{magnificat/magnificat-IIDsolemn}

  \begin{enumerate}[label=\textcolor{red}{\arabic*}]
    \setcounter{enumi}{2}
    \item Quia respéxit humilitátem \textit{an}\textit{cíl}\textit{læ} \textbf{su}æ:~* ecce enim ex hoc beátam me dicent omnes genera\textit{ti}\textbf{ó}nes.

    \item Quia fecit mihi \textit{ma}\textit{gna} \textit{qui} \textbf{pot}ens est:~* et sanctum no\textit{men} \textbf{e}jus.
    
    \item Et misericórdia ejus a progéni\textit{e} \textit{in} \textit{pro}\textbf{gé}nies~* timénti\textit{bus} \textbf{e}um.
    
    \item Fecit poténtiam in \textit{brá}\textit{chi}\textit{o} \textbf{su}o:~* dispérsit supérbos mente cor\textit{dis} \textbf{su}i.
    
    \item Depósuit pot\textit{én}\textit{tes} \textit{de} \textbf{se}de,~* et exaltá\textit{vit} \textbf{hú}miles.
    
    \item Esuriéntes \textit{im}\textit{plé}\textit{vit} \textbf{bo}nis:~* et dívites dimísit \textit{in}\textbf{á}nes.
    
    \item Suscépit Israël \textit{pú}\textit{e}\textit{rum} \textbf{su}um,~* recordátus misericórdi\textit{æ} \textbf{su}æ.
    
    \item Sicut locútus est \textit{ad} \textit{pa}\textit{tres} \textbf{nos}tros,~* Abraham et sémini ejus \textit{in} \textbf{sǽ}cula.
    
    \item Glória \textit{Pa}\textit{tri}, \textit{et} \textbf{Fí}lio,~* et Spirítu\textit{i} \textbf{Sanc}to.
    
    \item Sicut erat in princípio, \textit{et} \textit{nunc}, \textit{et} \textbf{sem}per,~* et in sǽcula sæculó\textit{rum}. \textbf{A}men.
  \end{enumerate}

  \grecommentary{\textit{Reprise de l'Antienne.}}
  \gabcsnippet{(f3) O(ehhg) A(h_f)do(gf)ná(ef)i,(f'_) (,) et(f) Dux(f) do(g)mus(hf) Is(g)ra(f)el,(e.) (;) qui(f) Mó(cf)y(f)si(f') (,) in(f) i(g)gne(f') flam(g)mae(f) ru(g)bi(f') (,) ap(f)pa(g')ru(f)í(g_[uh:l]h)sti,(kxiji/jkj.) (;) et(j) e(j')i(i) in(h) Si(h)na(h_g) le(h)gem(i') de(g)dí(hg)sti :(e.) (:) ve(eh)ni(g'_[oh:h]) (,) ad(h) red(hg)i(f)mén(gf~)dum(ef~) nos(f.) (,) in(ec~) brá(e)chi(f)o(h_g) ex(e)tén(gg)to.(f.)}

  \newpage

  \begin{center}
    \textcolor{red}{\large{Le 19 décembre}}
  \end{center}

  \gresetinitiallines{1}
  \greillumination{\initfamily\fontsize{11mm}{11mm}\selectfont O}
  \gregorioscore{antiennes/an--o_radix_jesse--solesmes_1961}
  \begin{center}
    \footnotesize{
      \textit{
        O Racine de Jessé, qui êtes comme l’étendard des peuples, (Isaïe XI, 10 et Romains XV, 12) devant qui les rois fermeront leur bouche, (Isaïe LII, 15) et dont les Nations imploreront le secours : venez nous délivrer, maintenant ne tardez plus. (Habacuc II, 3 et Hebreux X, 37.) 
      }
    }
  \end{center}
  \medskip

  \gresetinitiallines{0}
  \gregorioscore{magnificat/magnificat-IIDsolemn}

  \begin{enumerate}[label=\textcolor{red}{\arabic*}]
    \setcounter{enumi}{2}
    \item Quia respéxit humilitátem \textit{an}\textit{cíl}\textit{læ} \textbf{su}æ:~* ecce enim ex hoc beátam me dicent omnes genera\textit{ti}\textbf{ó}nes.

    \item Quia fecit mihi \textit{ma}\textit{gna} \textit{qui} \textbf{pot}ens est:~* et sanctum no\textit{men} \textbf{e}jus.
    
    \item Et misericórdia ejus a progéni\textit{e} \textit{in} \textit{pro}\textbf{gé}nies~* timénti\textit{bus} \textbf{e}um.
    
    \item Fecit poténtiam in \textit{brá}\textit{chi}\textit{o} \textbf{su}o:~* dispérsit supérbos mente cor\textit{dis} \textbf{su}i.
    
    \item Depósuit pot\textit{én}\textit{tes} \textit{de} \textbf{se}de,~* et exaltá\textit{vit} \textbf{hú}miles.
    
    \item Esuriéntes \textit{im}\textit{plé}\textit{vit} \textbf{bo}nis:~* et dívites dimísit \textit{in}\textbf{á}nes.
    
    \item Suscépit Israël \textit{pú}\textit{e}\textit{rum} \textbf{su}um,~* recordátus misericórdi\textit{æ} \textbf{su}æ.
    
    \item Sicut locútus est \textit{ad} \textit{pa}\textit{tres} \textbf{nos}tros,~* Abraham et sémini ejus \textit{in} \textbf{sǽ}cula.
    
    \item Glória \textit{Pa}\textit{tri}, \textit{et} \textbf{Fí}lio,~* et Spirítu\textit{i} \textbf{Sanc}to.
    
    \item Sicut erat in princípio, \textit{et} \textit{nunc}, \textit{et} \textbf{sem}per,~* et in sǽcula sæculó\textit{rum}. \textbf{A}men.
  \end{enumerate}

  \grecommentary{\textit{Reprise de l'Antienne.}}
  \gabcsnippet{(f3) O(ehhg) Ra(h_f)dix(gf) Jes(ef)se,(f'_) (,) qui(f) stas(fg) in(f) si(g)gnum(f) po(g_[uh:l]h)pu(fg)ló(f)rum,(e.) (;) su(f)per(cf~) quem(f) con(f)ti(f)né(fg)bunt(f) re(g')ges(f) os(g_[uh:l]h) su(kxiji)um,(jkj.) (;) quem(j_i~) gen(j)tes(i) de(h)pre(hg)ca(f!gwh)bún(h)tur :(e.) (:) ve(eh)ni(g'_[oh:h]) (,) ad(h) li(hg)be(f)rán(gf~)dum(ef~) nos,(fv_EC.) (;) jam(e_[uh:l]f) no(h_g)li(e) tar(gh)dá(f.)re.(f.) (::)}

  \newpage

  \begin{center}
    \textcolor{red}{\large{Le 20 décembre}}
  \end{center}

  \gresetinitiallines{1}
  \greillumination{\initfamily\fontsize{11mm}{11mm}\selectfont O}
  \gregorioscore{antiennes/an--o_clavis_david--solesmes_1961}
  \begin{center}
    \footnotesize{
      \textit{
        O Emmanuel (Isaïe VII, 14 et VIII, 8), notre Roi et notre Législateur (Isaïe XXXIII, 22), Attente des Nations (Genèse XLIX, 10) et leur Sauveur : venez nous sauver, Seigneur notre Dieu.  
      }
    }
  \end{center}
  \medskip

  \gresetinitiallines{0}
  \gregorioscore{magnificat/magnificat-IIDsolemn}

  \begin{enumerate}[label=\textcolor{red}{\arabic*}]
    \setcounter{enumi}{2}
    \item Quia respéxit humilitátem \textit{an}\textit{cíl}\textit{læ} \textbf{su}æ:~* ecce enim ex hoc beátam me dicent omnes genera\textit{ti}\textbf{ó}nes.

    \item Quia fecit mihi \textit{ma}\textit{gna} \textit{qui} \textbf{pot}ens est:~* et sanctum no\textit{men} \textbf{e}jus.
    
    \item Et misericórdia ejus a progéni\textit{e} \textit{in} \textit{pro}\textbf{gé}nies~* timénti\textit{bus} \textbf{e}um.
    
    \item Fecit poténtiam in \textit{brá}\textit{chi}\textit{o} \textbf{su}o:~* dispérsit supérbos mente cor\textit{dis} \textbf{su}i.
    
    \item Depósuit pot\textit{én}\textit{tes} \textit{de} \textbf{se}de,~* et exaltá\textit{vit} \textbf{hú}miles.
    
    \item Esuriéntes \textit{im}\textit{plé}\textit{vit} \textbf{bo}nis:~* et dívites dimísit \textit{in}\textbf{á}nes.
    
    \item Suscépit Israël \textit{pú}\textit{e}\textit{rum} \textbf{su}um,~* recordátus misericórdi\textit{æ} \textbf{su}æ.
    
    \item Sicut locútus est \textit{ad} \textit{pa}\textit{tres} \textbf{nos}tros,~* Abraham et sémini ejus \textit{in} \textbf{sǽ}cula.
    
    \item Glória \textit{Pa}\textit{tri}, \textit{et} \textbf{Fí}lio,~* et Spirítu\textit{i} \textbf{Sanc}to.
    
    \item Sicut erat in princípio, \textit{et} \textit{nunc}, \textit{et} \textbf{sem}per,~* et in sǽcula sæculó\textit{rum}. \textbf{A}men.
  \end{enumerate}

  \grecommentary{\textit{Reprise de l'Antienne.}}
  \gabcsnippet{(f3) O(ehhg) Ra(h_f)dix(gf) Jes(ef)se,(f'_) (,) qui(f) stas(fg) in(f) si(g)gnum(f) po(g_[uh:l]h)pu(fg)ló(f)rum,(e.) (;) su(f)per(cf~) quem(f) con(f)ti(f)né(fg)bunt(f) re(g')ges(f) os(g_[uh:l]h) su(kxiji)um,(jkj.) (;) quem(j_i~) gen(j)tes(i) de(h)pre(hg)ca(f!gwh)bún(h)tur :(e.) (:) ve(eh)ni(g'_[oh:h]) (,) ad(h) li(hg)be(f)rán(gf~)dum(ef~) nos,(fv_EC.) (;) jam(e_[uh:l]f) no(h_g)li(e) tar(gh)dá(f.)re.(f.) (::)}

  \newpage

  \begin{center}
    \textcolor{red}{\large{Le 21 décembre}}
  \end{center}

  \gresetinitiallines{1}
  \greillumination{\initfamily\fontsize{11mm}{11mm}\selectfont O}
  \gregorioscore{antiennes/an--o_oriens--solesmes_1961}
  \begin{center}
    \footnotesize{
      \textit{
        O Orient (Zacharie VI, 12), splendeur de la lumière éternelle (Hébreux I, 3), et soleil de justice (Malachie IV, 2) : venez et éclairez ceux qui sont assis dans les ténèbres et dans l’ombre de la mort (Isaïe IX, 2 & Luc I, 78-79). 
      }
    }
  \end{center}
  \medskip

  \gresetinitiallines{0}
  \gregorioscore{magnificat/magnificat-IIDsolemn}

  \begin{enumerate}[label=\textcolor{red}{\arabic*}]
    \setcounter{enumi}{2}
    \item Quia respéxit humilitátem \textit{an}\textit{cíl}\textit{læ} \textbf{su}æ:~* ecce enim ex hoc beátam me dicent omnes genera\textit{ti}\textbf{ó}nes.

    \item Quia fecit mihi \textit{ma}\textit{gna} \textit{qui} \textbf{pot}ens est:~* et sanctum no\textit{men} \textbf{e}jus.
    
    \item Et misericórdia ejus a progéni\textit{e} \textit{in} \textit{pro}\textbf{gé}nies~* timénti\textit{bus} \textbf{e}um.
    
    \item Fecit poténtiam in \textit{brá}\textit{chi}\textit{o} \textbf{su}o:~* dispérsit supérbos mente cor\textit{dis} \textbf{su}i.
    
    \item Depósuit pot\textit{én}\textit{tes} \textit{de} \textbf{se}de,~* et exaltá\textit{vit} \textbf{hú}miles.
    
    \item Esuriéntes \textit{im}\textit{plé}\textit{vit} \textbf{bo}nis:~* et dívites dimísit \textit{in}\textbf{á}nes.
    
    \item Suscépit Israël \textit{pú}\textit{e}\textit{rum} \textbf{su}um,~* recordátus misericórdi\textit{æ} \textbf{su}æ.
    
    \item Sicut locútus est \textit{ad} \textit{pa}\textit{tres} \textbf{nos}tros,~* Abraham et sémini ejus \textit{in} \textbf{sǽ}cula.
    
    \item Glória \textit{Pa}\textit{tri}, \textit{et} \textbf{Fí}lio,~* et Spirítu\textit{i} \textbf{Sanc}to.
    
    \item Sicut erat in princípio, \textit{et} \textit{nunc}, \textit{et} \textbf{sem}per,~* et in sǽcula sæculó\textit{rum}. \textbf{A}men.
  \end{enumerate}

  \grecommentary{\textit{Reprise de l'Antienne.}}
  \gabcsnippet{(f3) O(ehhg) O(h_f/gf)ri(ef)ens,(f.) <c>*</c>(,) splen(fg)dor(f) lu(f)cis(g_[uh:l]h) ae(fg)tér(f)nae,(e.) (;) et(f) sol(cf) ju(f)stí(g')ti(h)ae :(kxiji/jkj.) (:) ve(j)ni,(i) (,) et(hg) il(fh~)lú(h)mi(f_gF'E)na(e.) (;) se(h')dén(g)tes(h'_) in(f) té(gf)ne(ef)bris(fv_EC.) (;) et(e_[uh:l]f) um(h_g~)bra(e) mor(gg)tis.(f.) (::)}

  \newpage

  \begin{center}
    \textcolor{red}{\large{Le 22 décembre}}
  \end{center}

  \gresetinitiallines{1}
  \greillumination{\initfamily\fontsize{11mm}{11mm}\selectfont O}
  \gregorioscore{antiennes/an--o_rex_gentium--solesmes_1961}
  \begin{center}
    \footnotesize{
      \textit{
        O Roi des Nations, et objet de leurs désirs (Aggée II, 8), Pierre angulaire (Isaïe XXVIII, 16), qui réunissez en vous les deux peuples (Ephésiens II, 14) : venez et sauvez l’homme, que vous avez formé du limon (Genèse II, 7).  
      }
    }
  \end{center}
  \medskip

  \gresetinitiallines{0}
  \gregorioscore{magnificat/magnificat-IIDsolemn}

  \begin{enumerate}[label=\textcolor{red}{\arabic*}]
    \setcounter{enumi}{2}
    \item Quia respéxit humilitátem \textit{an}\textit{cíl}\textit{læ} \textbf{su}æ:~* ecce enim ex hoc beátam me dicent omnes genera\textit{ti}\textbf{ó}nes.

    \item Quia fecit mihi \textit{ma}\textit{gna} \textit{qui} \textbf{pot}ens est:~* et sanctum no\textit{men} \textbf{e}jus.
    
    \item Et misericórdia ejus a progéni\textit{e} \textit{in} \textit{pro}\textbf{gé}nies~* timénti\textit{bus} \textbf{e}um.
    
    \item Fecit poténtiam in \textit{brá}\textit{chi}\textit{o} \textbf{su}o:~* dispérsit supérbos mente cor\textit{dis} \textbf{su}i.
    
    \item Depósuit pot\textit{én}\textit{tes} \textit{de} \textbf{se}de,~* et exaltá\textit{vit} \textbf{hú}miles.
    
    \item Esuriéntes \textit{im}\textit{plé}\textit{vit} \textbf{bo}nis:~* et dívites dimísit \textit{in}\textbf{á}nes.
    
    \item Suscépit Israël \textit{pú}\textit{e}\textit{rum} \textbf{su}um,~* recordátus misericórdi\textit{æ} \textbf{su}æ.
    
    \item Sicut locútus est \textit{ad} \textit{pa}\textit{tres} \textbf{nos}tros,~* Abraham et sémini ejus \textit{in} \textbf{sǽ}cula.
    
    \item Glória \textit{Pa}\textit{tri}, \textit{et} \textbf{Fí}lio,~* et Spirítu\textit{i} \textbf{Sanc}to.
    
    \item Sicut erat in princípio, \textit{et} \textit{nunc}, \textit{et} \textbf{sem}per,~* et in sǽcula sæculó\textit{rum}. \textbf{A}men.
  \end{enumerate}

  \grecommentary{\textit{Reprise de l'Antienne.}}
  \gabcsnippet{(f3) O(ehhg) Rex(h_f) gén(gf~)ti(ef)um,(f'_) (,) et(f) de(f)si(f)de(f)rá(f)tus(g_[uh:l]h) e(fg)á(f)rum,(e.) (;) la(f)pís(cf)que(f) an(g')gu(f)lá(g_[uh:l]h)ris,(kxiji/jkj.) (;) qui(j) fa(j')cis(i) ú(h)tra(hg)que(f!gwh) u(h)num :(e.) (:) ve(eh)ni,(g'_[oh:h]) (,) et(h) sal(hg~)va(f) hó(gf)mi(ef)nem,(fv_EC.) (;) quem(e) de(f) li(h')mo(g) for(e)má(gg)sti.(f.) (::) }

  \newpage

  \begin{center}
    \textcolor{red}{\large{Le 23 décembre}}
  \end{center}

  \gresetinitiallines{1}
  \greillumination{\initfamily\fontsize{11mm}{11mm}\selectfont O}
  \gregorioscore{antiennes/an--o_emmanuel--solesmes_1961}
  \begin{center}
    \footnotesize{
      \textit{
        O Emmanuel (Isaïe VII, 14 & VIII, 8), * notre Roi et notre Législateur (Isaïe XXXIII, 22), Attente des Nations (Genèse XLIX, 10) et leur Sauveur : venez nous sauver, Seigneur notre Dieu.  
      }
    }
  \end{center}
  \medskip

  \gresetinitiallines{0}
  \gregorioscore{magnificat/magnificat-IIDsolemn}

  \begin{enumerate}[label=\textcolor{red}{\arabic*}]
    \setcounter{enumi}{2}
    \item Quia respéxit humilitátem \textit{an}\textit{cíl}\textit{læ} \textbf{su}æ:~* ecce enim ex hoc beátam me dicent omnes genera\textit{ti}\textbf{ó}nes.

    \item Quia fecit mihi \textit{ma}\textit{gna} \textit{qui} \textbf{pot}ens est:~* et sanctum no\textit{men} \textbf{e}jus.
    
    \item Et misericórdia ejus a progéni\textit{e} \textit{in} \textit{pro}\textbf{gé}nies~* timénti\textit{bus} \textbf{e}um.
    
    \item Fecit poténtiam in \textit{brá}\textit{chi}\textit{o} \textbf{su}o:~* dispérsit supérbos mente cor\textit{dis} \textbf{su}i.
    
    \item Depósuit pot\textit{én}\textit{tes} \textit{de} \textbf{se}de,~* et exaltá\textit{vit} \textbf{hú}miles.
    
    \item Esuriéntes \textit{im}\textit{plé}\textit{vit} \textbf{bo}nis:~* et dívites dimísit \textit{in}\textbf{á}nes.
    
    \item Suscépit Israël \textit{pú}\textit{e}\textit{rum} \textbf{su}um,~* recordátus misericórdi\textit{æ} \textbf{su}æ.
    
    \item Sicut locútus est \textit{ad} \textit{pa}\textit{tres} \textbf{nos}tros,~* Abraham et sémini ejus \textit{in} \textbf{sǽ}cula.
    
    \item Glória \textit{Pa}\textit{tri}, \textit{et} \textbf{Fí}lio,~* et Spirítu\textit{i} \textbf{Sanc}to.
    
    \item Sicut erat in princípio, \textit{et} \textit{nunc}, \textit{et} \textbf{sem}per,~* et in sǽcula sæculó\textit{rum}. \textbf{A}men.
  \end{enumerate}

  \grecommentary{\textit{Reprise de l'Antienne.}}
  \gabcsnippet{(f3) O(ehhg) Em(h_f)má(gf)nu(ef)el,(f'_) (,) Rex(g) et(f) lé(f)gi(g_[uh:l]h)fer(fg) no(f)ster,(e.) (;) ex(f)spe(cf)ctá(fg)ti(f)o(f) gén(g')ti(h)um,(kxiji/jkj.) (;) et(j') Sal(i)vá(h)tor(hg) e(f!gwh)á(h)rum :(e.) (:) ve(eh)ni(g'_[oh:h]) (,) ad(h) sal(f)ván(gf~)dum(ef~) nos(f.) (,) Dó(ec)mi(e)ne(f) De(h_g)us(e) no(gg)ster.(f.) (::)}

  \newpage
  \subfile{salut.tex}
\end{document}